%\documentclass[twoside,twocolumn]{article}
\documentclass[twoside,onecolumn,12pt]{article}

%\usepackage[sc]{mathpazo} % Use the Palatino font
\usepackage[T1]{fontenc} % Use 8-bit encoding that has 256 glyphs
\linespread{1.05} % Line spacing - Palatino needs more space between lines
%\usepackage{microtype} % Slightly tweak font spacing for aesthetics
\usepackage{graphicx}			% Inclusao de graficos

%\usepackage[latin1]{inputenc}
\usepackage[utf8]{inputenc}
%\inputencoding{utf8}
%\usepackage{ucs}


%\usepackage[english]{babel} % Language hyphenation and typographical rules
\usepackage[portuguese]{babel}

%\usepackage[hmarginratio=1:1,top=32mm,columnsep=20pt]{geometry} % Document margins
\usepackage{geometry} % Document margins
 \geometry{
 a4paper,
 total={170mm,257mm},
 left=20mm,
 top=20mm,
 }
 
\usepackage[hang, small,labelfont=bf,up,textfont=it,up]{caption} % Custom captions under/above floats in tables or figures
\usepackage{booktabs} % Horizontal rules in tables

\usepackage{setspace}

%\usepackage{lettrine} % The lettrine is the first enlarged letter at the beginning of the text

\usepackage{enumitem} % Customized lists
\setlist[itemize]{noitemsep} % Make itemize lists more compact

\usepackage{abstract} % Allows abstract customization
\renewcommand{\abstractnamefont}{\normalfont\bfseries} % Set the "Abstract" text to bold
\renewcommand{\abstracttextfont}{\normalfont\small\itshape} % Set the abstract itself to small italic text

\usepackage{titlesec} % Allows customization of titles
%\renewcommand\thesection{\Roman{section}} % Roman numerals for the sections
%\renewcommand\thesubsection{\roman{subsection}} % roman numerals for subsections
%\titleformat{\section}[block]{\large\scshape\centering}{\thesection.}{1em}{} % Change the look of the section titles
%\titleformat{\subsection}[block]{\large}{\thesubsection.}{1em}{} % Change the look of the section titles
\titleformat{\section}[block]{\large\bf\scshape}{\thesection.}{1em}{} % Change the look of the section titles
\titleformat{\subsection}[block]{\large\bf}{\thesubsection.}{1em}{} % Change the look of the section titles

\usepackage{fancyhdr} % Headers and footers
\pagestyle{fancy} % All pages have headers and footers
\fancyhead{} % Blank out the default header
\fancyfoot{} % Blank out the default footer
\fancyhead[C]{Econometria do Setor P\'{u}blico $\bullet$ Prof. Rafael Terra $\bullet$ 2\ord~sem. 2016} % Custom header text
\fancyfoot[RO,LE]{\thepage} % Custom footer text

% Incluir os pacotes booktabs e dcolumn, pois o comando esttab do Stata gera tabelas em latex que utilizam estes pacotes.
\usepackage{booktabs}
\usepackage{dcolumn}

\usepackage{titling} % Customizing the title section

\usepackage{hyperref} % For hyperlinks in the PDF

\usepackage{overpic}

\usepackage{indentfirst}

\usepackage{amsmath}

\usepackage{array}

\usepackage{multirow}

\usepackage{float}

\usepackage{color}				% Controle das cores
\usepackage[usenames,dvipsnames,svgnames,table]{xcolor}

\usepackage[brazilian,hyperpageref]{backref}	 % Paginas com as citacoes na bibl
%\usepackage[alf,versalete,abnt-etal-cite=5]{abntex2cite}	% Citacoes padrao ABNT
\usepackage[alf,abnt-etal-cite=5]{abntex2cite}	% Citacoes padrao ABNT

\usepackage[flushleft]{threeparttable}

%----------------------------------------------------------------------------------------
% Definicoes usadas para mostrar o codigo Stata
%----------------------------------------------------------------------------------------

\usepackage{listings}
\usepackage{fancyvrb}

\usepackage{accsupp}% http://ctan.org/pkg/accsupp
\newcommand{\emptyaccsupp}[1]{\BeginAccSupp{ActualText={}}#1\EndAccSupp{}}

\lstdefinestyle{numbers}{
    numbers=left, 
    stepnumber=1, 
    numberstyle=\tiny\emptyaccsupp,
    xleftmargin=2em,
}

\lstdefinestyle{nonumbers}{
    numbers=none,
}

\lstdefinelanguage{Stata}{
    % Left for users to add missing commands,
    % with possibility of choosing different style.
    morekeywords=,
    % Popular add-on commands
    morekeywords=[2]{cntrade, chinafin, 
                     wbopendata, spmap,
    },
    % System commands
    morekeywords=[3]{regress,xtreg,ivregress,ivreg2,summarize,display},
    % Keywords
    morekeywords=[4]{use,sort,merge,destring,keep,save,saveold,generate,recode,clear,gen,egen,by,bys,bysort,replace,drop},
    % Built-in functions
    morekeywords=[5]{forvalues, if, foreach, set, estat, hettest, local, return, list, sampsi, r, graph, twoway, kdensity, score, order, label, var, rnormal, runiform},
    morecomment=[l]{**},
    % morecomment=[l]{*},  // `*` maybe used as multiply operator. So use `//` as line comment.
    morecomment=[s]{/*}{*/},
    % morecomment=[s]{,}{//},
    % The following is used by macros, like `lags'.
    morecomment=[n][keywordstyle9]{`}{'},
    morestring=[b]",
    sensitive=true,
}

\lstdefinestyle{stata-plain}{
    % comment slanted and keywords bolded.
    language=Stata,
%    basicstyle=\setmonofont{Consolas}\footnotesize\ttfamily,
    basicstyle=\scriptsize\ttfamily,
    captionpos=t
}

\lstdefinestyle{stata-editor}{
    language=Stata,
    % size of the fonts for the code
%    basicstyle=\setmonofont{Consolas}\footnotesize\ttfamily,  
    basicstyle=\scriptsize\ttfamily,
    captionpos=t,
    % Color settings to match do-file editor style
    % Commands that are not included in the defination
    keywordstyle={\color{NavyBlue}},  
    % Popular add-on commands
    keywordstyle=[2]{\color{NavyBlue}},
    % System commands
    keywordstyle=[3]{\color{NavyBlue}},
    % Keywords
    keywordstyle=[4]{\color{NavyBlue}},
    % Built-in functions like rnormal
    keywordstyle=[5]{\color{Blue}},
    % Used by macros, i.e. `lags'  
    keywordstyle=[9]{\color{TealBlue}},
    stringstyle=\color{Maroon},
    commentstyle=\color{OliveGreen},
}

\definecolor{outputColor}{RGB}{235,235,235}

\definecolor{gainsboro}{RGB}{220,220,220}
\definecolor{lightgray}{RGB}{211,211,211}
\definecolor{silver}{RGB}{192,192,192}
\definecolor{darkgray}{RGB}{169,169,169}
\definecolor{gray}{RGB}{128,128,128}
\definecolor{dimgray}{RGB}{105,105,105}


\lstdefinestyle{OutputStyle}{ % 
  basicstyle=\scriptsize\ttfamily, % the size of the fonts that are used for the code
  numbers=none,                   % where to put the line-numbers
  backgroundcolor=\color{outputColor},  % choose the background color. You must add \usepackage{color}
  showspaces=false,               % show spaces adding particular underscores
  showstringspaces=false,         % underline spaces within strings
  showtabs=false,                 % show tabs within strings adding particular underscores
  frame=single,                   % adds a frame around the code
  rulecolor=\color{darkgray},        % if not set, the frame-color may be changed on line-breaks within not-black text (e.g. commens (green here))
  tabsize=2,                      % sets default tabsize to 2 spaces
  breaklines=true,                % sets automatic line breaking
  breakatwhitespace=false
} 

\lstnewenvironment{StataListing}{\lstset{language=Stata,style=stata-plain,style=numbers,showstringspaces=false,breaklines,frame=single}}{}

\lstnewenvironment{OutputListing}{\lstset{style=OutputStyle}}{}

\renewcommand{\lstlistingname}{C\'{o}digo-fonte}
\renewcommand*{\lstlistlistingname}{Lista de c\'{o}digos-fontes}


% Definicao do environment espec\'{i}ficos para as sa\'{i}das

\newcounter{outputctr}

\newenvironment{Output}[1][]{%      define a custom environment
   \par\noindent%         create a vertical offset to previous material
    \refstepcounter{outputctr}% increment the environment's counter
   {Sa\'{i}da \theoutputctr}% or \textbf, \textit, ...
    \if\relax\detokenize{~\textendash~#1}\relax
    \else
    ~\textendash~#1
    \fi
    \centering
   }{\par\bigskip}  %          create a vertical offset to following material
\numberwithin{outputctr}{section}

%----------------------------------------------------------------------------------------
%	TITLE SECTION
%----------------------------------------------------------------------------------------

\setlength{\droptitle}{-4\baselineskip} % Move the title up
\pretitle{\begin{center}\Huge\bfseries} % Article title formatting
\posttitle{\end{center}} % Article title closing formatting
\title{Impacto de Curto Prazo do Programa Mais M\'{e}dicos sobre o Estado de Saúde da População}
%\title{Efeitos do Programa Mais M\'{e}dicos}


\author{
      \textsc{Saulo Benchimol Bastos} \\[-1mm] % Your name
      \normalsize Universidade de Bras\'{i}lia -- UnB \\[-1mm] % Your institution
      \normalsize \href{mailto:sbenchimol@gmail.com}{sbenchimol@gmail.com} % Your email address
  \and % Uncomment if 2 authors are required, duplicate these 4 lines if more
      \textsc{Marcleition Ribeiro Morais} \\[-1mm] % Second author's name
      \normalsize Universidade de Bras\'{i}lia -- UnB \\[-1mm] % Second author's institution
      \normalsize \href{mailto:mrm@mail.uft.edu.br}{mrm@mail.uft.edu.br} % Second author's email address
}
%\date{\today} % Leave empty to omit a date
\date{9 de dezembro de 2016}
\renewcommand{\maketitlehookd}{%
\begin{abstract}
{\linespread{1.0}
\noindent O artigo parte do problema apontado pela literatura quanto a ineficácia do Programa Mais Médicos em propiciar melhoria no estado de saúde da população, e busca evidências de seus efeitos ainda no curto prazo. Para isso, faz uso da técnica de estimação por Diferenças em Diferenças, controlando por covariadas socioeconômicas, ambientais e relativas a serviços médicos. Tudo isso condicionado à existência de adequada correlação entre o estado de saúde da população e a propor\c{c}\~{a}o de nascidos vivos de m\~{a}es com sete ou mais consultas pr\'{e}-natal. Os resultados das estimações por efeitos fixos sugerem que o efeito do programa sobre a variável de interesse é positivo e bastante consistente com os controles. O impacto foi estimado em cerca de 1.8, em média, no modelo básico, controlando por aspectos socioeconsômicos, ambientais e no modelo com todos os controles. Em todos os casos, os resultados contribuem para validar a hipótese de que o programa já produziu efeito sobre a saúde da população em 2014, um ano após sua implantação.}

\end{abstract}
}
%O objetivo do artigo \'{e} avaliar se tanto as previs\~{o}es feitas por \citeonline{bibMendes2014} quanto
%o conte\'{u}do do relat\'{o}rio (TC 005.391/2014-8) do \citeonline{bibTCU2014} a respeito da inexist\^{e}ncia de impacto
%de curto prazo do Programa Mais M\'{e}dicos sobre a sa\'{u}de da popula\c{c}\~{a}o podem ser refutados a partir do uso das
%t\'{e}cnicas de avalia\c{c}\~{a}o de impacto, dentre elas \textit{Propensity Score Matching} e \textit{Diff-in-Diff}.
%\textcolor{red}{Adicionalmente, para contornar o problema apontado pelo TCU quanto \`{a} inexist\^{e}ncia de indicadores de avalia\c{c}\~{a}o
%da sa\'{u}de da popula\c{c}\~{a}o, busca-se evid\^{e}ncias da qualidade dos instrumentos utilizados por Mendes, al\'{e}m de
%apresentar alternativas a sua representa\c{c}\~{a}o}. Os resultados s\~{a}o conclusivos...}


% ---
% Configuracoes de aparencia do PDF final

% alterando o aspecto da cor azul
\definecolor{blue}{RGB}{41,5,195}

% informacoes do PDF
\makeatletter
\hypersetup{
     	%pagebackref=false,
	pdftitle={\@title}, 
	pdfauthor={Saulo Benchimol Bastos e Marcleiton Morais},
	pdfsubject={assunto? econometria!},
	pdfcreator={Saulo Benchimol Bastos e Marcleition Morais},
	pdfkeywords={econometria}{avalia\c{c}\~{a}o de programas}{trabalho acad\^{e}mico}, 
	colorlinks=true,       	% false: boxed links; true: colored links
    	linkcolor=blue,         % color of internal links
    	citecolor=blue,        	% color of links to bibliography
    	filecolor=magenta,      % color of file links
	urlcolor=blue,
	bookmarksdepth=4
}
\makeatother

\makeatletter
\newcommand\urlfootnote@[1]{\footnote{\url@{#1}}}
\DeclareRobustCommand{\urlfootnote}{\hyper@normalise\urlfootnote@}
\makeatother

% ---
% Configuracoes do pacote backref
% Usado sem a opcao hyperpageref de backref
\renewcommand{\backrefpagesname}{~}
% Texto padrao antes do numero das paginas
\renewcommand{\backref}{}
% Define os textos da citacao
\renewcommand*{\backrefalt}[4]{~}%
% ---

\citeoption{abnt-and-type=e}
\citeoption{abnt-etal-cite=5}

%----------------------------------------------------------------------------------------

\begin{document}

% Print the title
\maketitle

%----------------------------------------------------------------------------------------
%	ARTICLE CONTENTS
%----------------------------------------------------------------------------------------

\setstretch{1.5}

\section{Introdu\c{c}\~{a}o \label{secIntroducao}}

%\lettrine[nindent=0em,lines=3]{L}

O Programa Mais Médicos lançado em meados de 2013 pelo Governo Federal é mais uma iniciativa de promover melhoria do atendimento aos usu\'{a}rios do Sistema \'{U}nico de Sa\'{u}de (SUS), corrigir o problema de distorção de oferta de médicos no país, e ampliar os investimentos para constru\c{c}\~{a}o, reforma e amplia\c{c}\~{a}o de Unidades B\'{a}sicas de Sa\'{u}de (UBS). Além disso, o programa busca ofertar novas vagas de gradua\c{c}\~{a}o, e resid\^{e}ncia m\'{e}dica para qualificar a forma\c{c}\~{a}o de pessoal na área.

Suas ações são orientadas por três eixos norteadores: infraestrutura, educação e provimento emergencial. Quanto à infraestrutura, busca-se ampliar e qualificar as instala\c{c}\~{o}es das UBS. Em termos educacionais, a intenção é ampliar o quantitativo de cursos e vagas na gradua\c{c}\~{a}o e na resid\^{e}ncia m\'{e}dica, com \^{e}nfase para as regi\~{o}es com menor rela\c{c}\~{a}o de vagas e m\'{e}dicos por habitante, al\'{e}m de promover mudan\c{c}as e reorienta\c{c}\~{a}o na forma\c{c}\~{a}o m\'{e}dica. Por sua vez, o provimento emergencial provisiona m\'{e}dicos para as \'{a}reas e munic\'{i}pios priorit\'{a}rios.

A implantação do programa porém incorreu quase que unicamente sobre este o eixo emergencial. Por isso, desde então, o programa vem sendo alvo de diversas pesquisas com o intuito de averiguar se sua intervenção produziria impacto sobre o estado de saúde da população. Para além da mudança de indicadores diretamente relacionados ao programa, dois estudos recentes apontaram para a inexistência de efeito real sobre a saúde da população, colocando dúvidas sobre a efetividade do gasto público.

\citeonline{bibMendes2014} usou dados em painel de 2008 a 2010 dos municípios brasileiros envolvendo variáveis de ordem sócioeconômica e do estado de saúde da população, para sugerir que o aumento do número de médicos, isoladamente,  não seria capaz de gerar melhoras na taxa de mortalidade infantil. No mesmo ano, o programa foi objeto de avaliação do Tribunal de Contas da União, \citeonline{bibTCU2014}, por meio do Relat\'{o}rio de Auditoria Operacional do Programa Mais M\'{e}dicos (TC 005.391/2014-8). Apesar da avaliação apontar para efeitos positivos sobre os indicadores de atendimento do SUS, o relatório sugere que seus efeitos não seriam significativos para o estado de saúde da população no curto prazo.

É notório que o programa possui diversas frentes de atua\c{c}\~{a}o, algumas de longo, e outras de curto prazo. Ademais, ele incide sobre uma realidade adversa da saúde pública brasileira, já que, em muitos municípios observa-se uma completa escassez de serviços médicos. Nesse sentido, considera-se razo\'{a}vel esperar que o programa produza efeito já no curto prazo.

Esta pesquisa tem como objetivo principal de estimar o efeito do Programa Mais M\'{e}dicos sobre o estado de sa\'{u}de da popula\c{c}\~{a}o no curto prazo, fazendo uso da técnicas de avaliação de impacto Diferenças em Diferenças. Além disso, busca-se evidências para a existência de efeito proporcionalmente maior em municípios com baixa renda.

%\textcolor{red}{Adicionalmente, para contornar o problema apontado pelo TCU quanto \`{a} inexist\^{e}ncia de indicadores de avalia\c{c}\~{a}o da sa\'{u}de da popula\c{c}\~{a}o, busca-se evid\^{e}ncias da qualidade dos instrumentos utilizados por Mendes, al\'{e}m de apresentar alternativas a sua representa\c{c}\~{a}o}.

%------------------------------------------------

\section{Referencial Te\'{o}rico \label{secReferencialTeorico}}

%\subsection{Programa Mais M\'{e}dicos}
%
%O Programa Mais M\'{e}dicos \cite{bibProgramaMaisMedicos2015,bibConhecaPMM2016} \'{e}
%uma iniciativa do Governo Federal que tem como objetivos:
%\begin{itemize}
%  \item Melhorar o atendimento aos usu\'{a}rios do Sistema \'{U}nico de Sa\'{u}de (SUS);
%  \item Levar mais m\'{e}dicos para regi\~{o}es onde h\'{a} escassez ou aus\^{e}ncia desses profissionais;
%  \item Ampliar os investimentos para constru\c{c}\~{a}o, reforma e amplia\c{c}\~{a}o de Unidades B\'{a}sicas de Sa\'{u}de (UBS);
%  \item Ofertar novas vagas de gradua\c{c}\~{a}o, e resid\^{e}ncia m\'{e}dica para qualificar a forma\c{c}\~{a}o desses profissionais.
%\end{itemize}
%
%Ele est\'{a} calcado em algumas frentes ou eixos:
%\begin{itemize}
%  \item \textbf{Infraestrura}: ampliar e qualificar as instala\c{c}\~{o}es das UBS;
%
%  \item \textbf{Educa\c{c}\~{a}o}: ampliar o quantitativo de cursos e vagas na gradua\c{c}\~{a}o e na
%    resid\^{e}ncia m\'{e}dica, com \^{e}nfase para as regi\~{o}es com menor rela\c{c}\~{a}o de vagas e
%    m\'{e}dicos por habitante, al\'{e}m de promover mudan\c{c}as e reorienta\c{c}\~{a}o na forma\c{c}\~{a}o
%    m\'{e}dica;
%
%  \item \textbf{Provimento emergencial}: prover - ou provisionar - m\'{e}dicos para as \'{a}reas e
%  munic\'{i}pios priorit\'{a}rios.
%\end{itemize}
%
%Este trabalho se insere no contexto do terceiro eixo.
%
%%------------------------------------------------

\subsection{Modelagem em saúde humana}\label{secModelosSaude}
%\subsection{Modelos propostos por \citeonline{bibGrossman1972} \label{secModelosGrossman}

Ao considerar a sa\'{u}de um tipo de estoque de capital humano que se deprecia ao longo do tempo a partir de um estoque inicial, \citeonline{bibGrossman1972} prop\~{o}e a estima\c{c}\~{a}o da demanda por saúde a partir do estudo da resposta do estoque de sa\'{u}de, do investimento bruto e outros determinantes de sa\'{u}de frente a altera\c{c}\~{o}es em vari\'{a}veis ex\'{o}genas. Para isso, ele prop\~{o}e dois casos especiais: o modelo de Investimento Puro e o de Consumo Puro. O autor recomenda o uso do modelo de investimento por que ele gera predi\c{c}\~{o}es mais consistentes para an\'{a}lises simples e os aspectos de consumo da demanda de sa\'{u}de podem ser incorporados em estima\c{c}\~{o}es emp\'{i}ricas sem muita perda de generalidade. A forma estrutural reduzida do modelo de Investimento Puro \'{e} dada pelo sistema abaixo:
\begin{subequations}
\begin{align}\label{eq_Demandasaude}
        \ln{\textbf{H}} = \alpha + \beta_W{\ln{\textbf{W}}} + \beta_P \ln{\textbf{P}} + \beta_E{\textbf{E}} + \beta_t{\textbf{t}} + \textbf{u},\\
        \ln{\textbf{M}} = \theta + \beta_{WM}{\ln{\textbf{W}}} +\beta_{PM}\ln{\textbf{P}} + \beta_{EM}{\textbf{E}} + \beta_{tM}{\textbf{t}} + \textbf{e}.\label{eq_Servicos}
\end{align}
\end{subequations}

Considerando a fun\c{c}\~{a}o de produ\c{c}\~{a}o (\ref{eq_Demandasaude}), o estoque de sa\'{u}de $\textbf{H}$ \'{e} end\'{o}geno em rela\c{c}\~{a}o ao tempo, $\textbf{W}$ \'{e} a taxa de sal\'{a}rio, $\textbf{E}$ \'{e} o estoque de capital humano que n\~{a}o capital sa\'{u}de humana e \'{e} ex\'{o}geno, $\textbf{P}$ \'{e} o vetor de pre\c{c}os dos bens considerados insumos da fun\c{c}\~{a}o de produ\c{c}\~{a}o de investimentos em sa\'{u}de, e $\textbf{u}$ é o termo de perturba\c{c}\~{a}o. O modelo prediz os seguintes resultados de interesse: $\beta_W>0$, $\beta_P<0$, $\beta_E>0$ e $\beta_t<0$. Servi\c{c}os m\'{e}dicos, $\textbf{M}$, s\~{a}o determinados pela equa\c{c}\~{a}o (\ref{eq_Servicos}), sendo que sua fun\c{c}\~{a}o de produ\c{c}\~{a}o secund\'{a}rio na presente an\'{a}lise. As vari\'{a}veis ex\'{o}genas s\~{a}o a taxa de sal\'{a}rio, pre\c{c}o dos cuidados m\'{e}dicos, estoque de capital humano, idade, e vari\'{a}vel latente taxa de deprecia\c{c}\~{a}o da sa\'{u}de no per\'{i}odo inicial.

Alguns fatos estilizados sobre o modelo \'{e} que a demanda por sa\'{u}de tem uma elasticidade renda zero, mas uma elasticidade renda positiva no modelo de consumo puro, tal que sa\'{u}de pode ser considerada um bem superior. Al\'{e}m disso, escolaridade formal conclu\'{i}da tem efeito positivo e \'{e} o mais importante determinante de estoque de capital humano, incluindo sa\'{u}de. Por fim, um acr\'{e}scimo na idade reduz a sa\'{u}de e aumenta os gastos m\'{e}dicos.

Essa estrutura deve considerar por\'{e}m outros mercados que podem determinar investimentos em sa\'{u}de, logo relacionados com $\textbf{P}$, como dietas, recrea\c{c}\~{a}o, consumo de cigarro e consumo excessivo de \'{a}lcool. Diversos fatores de ambiente tamb\'{e}m est\~{a}o diretamente relacionadas com a sa\'{u}de humana, como a condi\c{c}\~{a}o de habita\c{c}\~{a}o, saneamento b\'{a}sico, condi\c{c}\~{a}o de trabalho, entre outros.

A interação entre mercados bem como a dificuldade de se separar efeitos de oferta sobre a demanda de saúde é ainda mais difícil quando se considera o problema de ordem agregada. O fato dos serviços de saúde serem geralmente prestados em função da "necessidade" e não da "disposição a pagar", a prestação pública da maioria dos serviços de saúde, a hipótese subjacente de preferências homogêneas, juntamente com várias formas de não-racionamento de preços (como os tempos de espera) e registro não rigoroso dos dados são fatores que dificultam a estimação da demanda agregada [\citeonline{culyer2000handbook}].

Considerando-se despesas médicas como investimento em saúde, deve haver certo paralelismo entre a abordagem de \citeonline{bibGrossman1972} e a análise em dados agregados recente. Esta literatura trata do problema de damanda agregada por saúde considerando sua estimação a partir das despesas com cuidados médicos per capita entre países. Nesse sentido, \citeonline{gerdtham1998determinants} modelam o problema a partida da seguinte função de despesas: 

\begin{equation}
HE_{it} = b_0 + \sum_{j = 1}^{14}{ b_{j} X_{itj} } + \sum_{j = 1}^{10}{ c_{j} Z_{itj} } + \sum_{i}^{22}{\mu_{0i}b_{i}} + \sum_{t}^{22}{\eta_{0t}b_{t}} + e_{it}
\end{equation}
O modelo é log-linear, pois as variáveis contínuas são dispostas em logaritmo natural. Nele, $HE_{it}$ são as despesas com saúde; $X_{itj}$ inclui variáveis representativas de renda em termos per capita, estrutura e transição demográfica, condições de trabalho, hábito permanente com efeito sobre a saúde, despesas de pacientes com cuidados médicos, gastos públicos com o sistema de saúde, morbidade, número de médicos por 1000 habitantes, modalidades de serviços médicos; e $Z_{itj}$ é um conjunto de \textit{dummies} sobre reembolsos públicos, integração do sistema público saúde, limites orçamentários, sistema de manutenção, pagamentos diretos de pacientes, salários, entre outros. O termo $\sum_{i}^{22}{\mu_{0i}b_{i}} + \sum_{t}^{22}{\eta_{0t}b_{t}}$ tem o objetivo de corrigir problemas de multicolinearidade das unidades no tempo. Suas estimativas reforçam vários resultados da literatura sobre o impacto de variáveis institucionais e não institucionais sobre despesas com saúde.

A pesquisa atual considera o nível de agregação municipal de tal modo que o modelo de demanda agregada assume a seguinte forma geral:

\begin{equation}
H_{it} = \beta^{'}X_{it} + \rho^{'}M_{it} + \eta^{'}Z_{it} + v_i + e_{it}
\end{equation}

A forma funcional também é log-linear já que algumas variáveis são tratadas em logaritmo natural. Aqui, $H$ é o indicador de estado de saúde em nível municipal, $X$ representa o conjunto das variáveis socioeconômicas, $M$  corresponde aos serviços médicos, incluindo infraestrutura de atendimento geral e despesas próprias com saúde, e $Z$ as variáveis ambientais como fornecimento de água tratada, coleta de lixo e esgoto. Por fim, $v_i$ representa o efeito das características individuais do i-ésimo município que são invariantes no tempo.


\section{Metodologia \label{secMetodologia}}

%------------------------------------------------

\subsection{Estimador de diferen\c{c}as em diferen\c{c}as \label{secEstimadorDiffDiff}}

A estimação do modelo geral tem como base a análise de avaliação de impacto dada por meio do estimador de diferenças em diferenças. \citeonline{bibMicroeconometricsCameron} e \cite{bibAngrist2009mostly} mostram que,
para o modelo:
\begin{equation}
 y_{ist} = \alpha + \gamma_s + \lambda_t + \beta D_{st} + \varepsilon_{ist}
 \label{eq_diff_in_diff}
\end{equation}
\noindent onde $i$ corresponde ao indiv\'{i}duo, $s$ ao estado, e $t$ ao per\'{i}odo,
$\lambda_t$ \'{e} um efeito-fixo temporal, $\gamma_s$ est\'{a} relacionado ao estado,
$D_{st}$ \'{e} a vari\'{a}vel dummy de tratamento dada por:

\begin{equation}
 D_{st} = \left\{
 \begin{array}{rl}
  1 & \text{se houve tratamento dado pelo estado} ~ s ~ \text{no per\'{i}odo} ~ t \\
  0 & \text{caso contr\'{a}rio}
 \end{array}
 \right.
\end{equation}

\noindent o estimador de diferen\c{c}as em diferen\c{c}as \'{e} dado por:
\begin{equation}
\begin{array}{rcl}
 \hat{\beta} & = & E[\Delta y_{ist} | s = \text{T}] - E[\Delta y_{ist} | s = \text{NT}] \\ 
 ~           & = & (E[y_{ist} | s = \text{T}, t = t_2] - E[y_{ist} | s = \text{T}, t = t_1]) \\
 ~           & ~ & - (E[y_{ist} | s = \text{NT}, t = t_2] - E[y_{ist} | s = \text{NT}, t = t_1])
\end{array}
\end{equation}
\noindent sendo que $s=T$ corresponde ao estado tratado, $s=NT$ corresponde ao estado n\~{a}o-tratado,
$t=t_1$ \'{e} um per\'{i}odo pr\'{e}-tratamento e $t=t_2$ \'{e} um per\'{i}odo p\'{o}s-tratamento.

%------------------------------------------------

\subsubsection{Teste de tend\^{e}ncia no tempo \label{secTendenciaNoTempo}}

\citeonline{bibMicroeconometricsCameron} afirmam que,
quando se utiliza o estimador de diferen\c{c}as em diferen\c{c}as,
assume-se que os efeitos temporais s\~{a}o comuns entre os indiv\'{i}duos tratados e de controle.
Em outras palavras, que a tend\^{e}ncia \'{e} a mesma para os grupos de controle e de tratamento 
na aus\^{e}ncia de tratamento.

\citeonline{bibPischke2007} e \citeonline{bibAngrist2009mostly} prop\~{o}em uma forma de testar esta hip\'{o}tese,
que se baseia na inclus\~{a}o de avan\c{c}os e atrasos no efeito do tratamento:
\begin{equation}
 y_{ist} = \gamma_s + \lambda_t + \sum_{j = -m}^{q}{ \beta_{j} D_{s, k_s + j} } + X_{ist} \delta + \varepsilon_{ist}
 \label{eq_TendenciaNoTempo}
\end{equation}
\noindent onde $k_s$ \'{e} o tempo em que ocorreu o tratamento no estado $s$,
$D_{s,t}$ s\~{a}o as ``vari\'{a}veis de tratamento'', sendo que existem
``m'' atrasos e ``q'' avan\c{c}os no efeito do tratamento,
$\beta_{j}$ \'{e} o coeficiente do $j$-\'{e}simo avan\c{c}o ou atraso.
O teste do estimador de diferen\c{c}as em diferen\c{c}as \'{e} verificar se $\beta_{j} = 0$, $\forall j < 0$,
que essencialmente \'{e} o teste de causalidade de Granger \cite{bibAppliedEconometricTimeSeries2014}.
$\beta_{j}$ n\~{a}o precisa ser igual para $j \geq 0$, pois o efeito do tratamento pode acumular ao longo do tempo.

%------------------------------------------------

\subsection{Banco de dados}

Os dados utilizados est\~{a}o em n\'{i}vel municipal e correspondem a uma parcela das informa\c{c}\~{o}es do estado de sa\'{u}de da popula\c{c}\~{a}o brasileira fornecida pelo DATASUS, Departamento de Inform\'{a}tica do Sistema \'{U}nico de Sa\'{u}de (SUS). S\~{a}o indicadores que abrangem desde estat\'{i}sticas vitais a determinantes populacionais e da situa\c{c}\~{a}o sanit\'{a}ria. Grande parte est\'{a} acess\'{i}vel por meio do aplicativo TABNET e agrupados em diversos sistemas, como o Sistema de Informa\c{c}\~{a}o da Aten\c{c}\~{a}o B\'{a}sica (Siab), Sistema de Informa\c{c}\~{a}o Hospitalar (SIH), Informa\c{c}\~{o}es sobre Nascidos Vivos (Sinasc), Sistema de Informa\c{c}\~{a}o sobre Mortalidade (SIM), Sistema de Informa\c{c}\~{o}es sobre Or\c{c}amentos P\'{u}blicos em Sa\'{u}de (SIOPS) e  no Cadastro Nacional de Estabelecimentos de Sa\'{u}de (CNES), entre outros. Vale ressaltar que algumas informa\c{c}\~{o}es foram obtidas junto ao Instituto Brasileiro de Geografia e Estat\'{i}stica (IBGE) e ao Instituto Nacional de Estudos e Pesquisas Educacionais An\'{i}sio Teixeira (INEP). Um resumo dessas bases, bem como seus endere\c{c}os eletr\^{o}nicos, \'{e} mostrado na Tabela \ref{tabBDs} (Ap\^{e}ndice \ref{secApendice_BDs}).

%\subsection{Descri\c{c}\~{a}o das vari\'{a}veis}
\subsubsection{Vari\'{a}veis dependentes}

Considerando que a sa\'{u}de humana \'{e} fun\c{c}\~{a}o de m\'{u}ltiplos fatores individuais e de ambiente que n\~{a}o os cuidados m\'{e}dicos, al\'{e}m de fatores din\^{a}micos de tempo, de onde pode-se originar endogenia, avaliou-se como fundamental a an\'{a}lise da qualidade da \textit{proxy} utilizada para representar a sa\'{u}de humana. As vari\'{a}veis de interesse para esta an\'{a}lise est\~{a}o dispostas na Tabela \ref{tabDVarDependente} e podem ser agrupadas em dois grupos, as vari\'{a}veis de impacto indireto compat\'{i}veis com o longo prazo e as vari\'{a}veis de impacto direto com efeito de curto prazo.

\begin{table}[!htb]
\caption{Descri\c{c}\~{a}o das Vari\'{a}veis Dependentes \label{tabDVarDependente}}
%\centering
{\footnotesize
\begin{tabular}{cp{7cm}p{3cm}p{2cm}}
\toprule\toprule
\textbf{Vari\'{a}vel} & \textbf{Descri\c{c}\~{a}o} & \textbf{Sinal Esperado} & \textbf{Fonte}\\
\midrule
\midrule
$txmortalidade$ & Taxa de Mortalidade Infantil \begin{equation*} txmortalidade=\frac{obitoinfatil}{nvpn}*1000 \label{Eqtxmortalidade}\end{equation*}  & - & SIM/Sinasc \\
\midrule
$txmortalidadeadulta$ & Taxa de Mortalidade Adulta \begin{equation*} txmortalidadeadulta=\frac{obitosdoenca}{obitostotal}*1000 \label{Eqtxmortalidadeadulta}\end{equation*}& - & SIM\\
\midrule
$txprenatal$ & Propor\c{c}\~{a}o de nascidos vivos de m\~{a}es com sete ou mais consultas de pr\'{e}-natal \begin{equation*} txprenatal=\frac{nv7oumaisconsult}{nvpn}*100 \label{Eqtxprenatal}\end{equation*} & + & Sinasc \\[8mm]
\midrule
$acompbf$ & Cobertura de acompanhamento das condicionalidades de Sa\'{u}de do Programa Bolsa Fam\'{i}lia \begin{equation*} acompbf=\frac{familiaspbfacompat\_basica}{familiaspbf}*100 \label{Eqacompbf}\end{equation*}& + & PBF\\
\midrule
$icsab$ & Propor\c{c}\~{a}o de interna\c{c}\~{o}es por condi\c{c}\~{o}es sens\'{i}veis \`{a} Aten\c{c}\~{a}o B\'{a}sica \begin{equation*} icsab=\frac{interncond\_sensiv\_at\_basica}{internacoesclinicasnosus}*100 \label{Eq_icsab}\end{equation*}& -+ & SIH\\
\midrule
$consforaarea$ & Consultas m\'{e}dicas de pessoas residentes em \'{a}reas fora da abrang\^{e}ncia da equipe (PSF) & + & Siab\\
\bottomrule
\end{tabular}
\begin{tablenotes}
\item \tiny{Fonte: Elaboração dos autores.}
\end{tablenotes}
}
\end{table}

Do conjunto das vari\'{a}veis com efeito esperado mais de longo prazo, a Taxa de Mortalidade Infantil ($txmortalidade$) \'{e} definida como a raz\~{a}o entre o n\'{u}mero de \'{o}bitos em menores de 1 ano de idade em um determinado local de resid\^{e}ncia e ano e o n\'{u}mero de nascidos vivos residentes nesse mesmo local e ano, em unidades de mil habitantes. Seus dados t\^{e}m origem no SIM e no Sinasc, respectivamente. Considerou-se tamb\'{e}m a raz\~{a}o dos \'{o}bitos prematuros (30 a 69 anos) pelo conjunto das 4 principais doen\c{c}as cr\^{o}nicas n\~{a}o transmiss\'{i}veis (doen\c{c}as do aparelho circulat\'{o}rio, c\^{a}ncer, diabetes e doen\c{c}as respirat\'{o}rias cr\^{o}nicas) com o total de \'{o}bitos n\~{a}o fetais, ambos do SIM, de onde se originou a vari\'{a}vel Taxa de Mortalidade Adulta ($txmortalidadeadulta$). A inten\c{c}\~{a}o \'{e} inferir sobre h\'{a}bitos permanentes da popula\c{c}\~{a}o em rela\c{c}\~{a}o a condi\c{c}\~{a}o de sa\'{u}de, como o consumo de \'{a}lcool e tabaco.

Na tentativa de caracterizar o impacto direto dos cuidados m\'{e}dicos, utilizou-se a Propor\c{c}\~{a}o de Nascidos Vivos de M\~{a}es com Sete ou Mais Consultas de Pr\'{e}-natal no conjunto de todos os nascidos vivos ($txprenatal$), de fonte Sinasc.
% A fim de verificar se haveria uma invers\~{a}o ou aumento do efeito dentro das classes de atendimento menores, prop\^{o}s-se uma redefini\c{c}\~{a}o desse indicador a partir da propor\c{c}\~{a}o de nascidos vivos com menos de sete consultas, que originou as vari\'{a}veis $txprenatal0$, $txprenatal3$ e $txprenatal6$ \textcolor{red}{[falta definir essas vari\'{a}veis...]}.
Uma suspeita levantada \'{e} de que o atendimento m\'{e}dico n\~{a}o emergencial prestado \`{a} popula\c{c}\~{a}o mais pobre pode ser afetado pela presen\c{c}a de risco moral, \`{a} medida que o p\'{u}blico alvo n\~{a}o considere importante o usufruto de um servi\c{c}o t\~{a}o especializado, ou ao menos na frequ\^{e}ncia requerida, como preconizado por \citeonline{bibTheProductionOfHealth1972}.

Esse agrupamento disp\~{o}e ainda do indicador Cobertura de Acompanhamento das Condicionalidades de Sa\'{u}de do Programa Bolsa Fam\'{i}lia ($acompbf$), que \'{e} obtido pela raz\~{a}o entre o n\'{u}mero de fam\'{i}lias benefici\'{a}rias do Programa Bolsa Fam\'{i}lia (PBF) com perfil sa\'{u}de acompanhadas pela aten\c{c}\~{a}o b\'{a}sica na \'{u}ltima vig\^{e}ncia do ano e o  n\'{u}mero total de fam\'{i}lias beneci\'{a}rias do PBF com perfil sa\'{u}de na \'{u}ltima vig\^{e}ncia do ano, dados do Sistema de Gest\~{a}o do Acompanhamento das Condicionalidades de Sa\'{u}de do PBF. A inten\c{c}\~{a}o \'{e} monitorar o efeito do programa sobre a popula\c{c}\~{a}o em situa\c{c}\~{a}o de extrema pobreza, respondendo assim a um dos principais objetivos do programa de ampliar a oferta de servi\c{c}os m\'{e}dicos em munic\'{i}pios com essa caracter\'{i}stica. A vari\'{a}vel Propor\c{c}\~{a}o de Interna\c{c}\~{o}es por Condi\c{c}\~{o}es Sens\'{i}veis \`{a} Aten\c{c}\~{a}o B\'{a}sica ($icsab$) definida pela raz\~{a}o entre o n\'{u}mero de interna\c{c}\~{o}es por causas sens\'{i}veis selecionadas \`{a} aten\c{c}\~{a}o b\'{a}sica em determinado local e per\'{i}odo e o total de interna\c{c}\~{o}es cl\'{i}nicas em determinado local e per\'{i}odo, do SIH, e a vari\'{a}vel ($consforaarea$) correspondente ao total de consultas m\'{e}dicas de pessoas residentes em \'{a}reas fora da abrang\^{e}ncia da equipe do Programa Sa\'{u}de da Fam\'{i}lia (PSF), tamb\'{e}m atendem a esse prop\'{o}sito.

\subsubsection{Vari\'{a}veis independentes}


As vari\'{a}veis independentes utilizadas nas regress\~{o}es est\~{a}o detalhadas
na Tabela \ref{tabVariaveis} (Ap\^{e}ndice \ref{secApendice_Variaveis}). A dimensão socioeconômica abrange as medidas de escolaridade, aqui expressas pelas vari\'{a}veis taxa de abandono escolar do ensino médio ($abandono\_em$) e taxa de distor\c{c}\~{a}o idade-s\'{e}rie do aluno no ensino fundamental ($distorcao\_ef$), dados do INEP; e renda real per capita dada pelo PIB real per capita ($pibrealpc$). Vale ressaltar que os dados de PIB municipal são do IBGE e est\~{a}o dispon\'{i}veis somente at\'{e} o ano 2013. Para estimar os dados de 2014, foi feita a seguinte regress\~{a}o em nível de munic\'{i}pio:
\begin{equation}
 PIB(ano) = \alpha_0 + \alpha_1 \cdot ano
 \label{eq_estimativa_pib_municipal_1}
\end{equation}

% \begin{equation}
%  PIB(ano) = \alpha_0 + \alpha_1 \cdot PIB(ano-1)
%  \label{eq_estimativa_pib_municipal_2}
% \end{equation}

\noindent onde o valor do PIB para o ano de 2014 assume a forma $PIB(2014) = \alpha_0 + \alpha_1 \cdot 2014$.
% pela eq. (\ref{eq_estimativa_pib_municipal_1})
% e $PIB(2014) = \alpha_0 + \alpha_1 \cdot PIB(2013)$ pela eq. (\ref{eq_estimativa_pib_municipal_2}).

As vari\'{a}veis referentes \`{a}s condi\c{c}\~{o}es de saneamento s\~{a}o ($aguatratada$) e ($coletalixo$). Agua tratada corresponde ao total de domicílios com acesso a \'{a}gua tratada pela rede p\'{u}blica ou que utilizam \'{a}gua filtrada  ou \'{a}gua fervida ou, ainda, \'{a}gua clorada. Esse dados são do Siab, e correspondem aos registros realizados pro agentes de saúde sobre a abrangência dessas indicadores. Por fim, coleta de lixo é o total de domicílios com acesso aos serviços de coleta de lixo seja pública ou privada.

As \textit{proxies} para serviços médicos são despesa anual com saúde oriundas de recursos próprios do município ($despesaspsreal$), total de equipamentos existentes em uso pelo SUS tendo como base o m\^{e}s de maio do ano ($equipamentos$), dados do SIOPS, e o n\'{u}mero de equipes de saúde da família (eSF) aptas para o munic\'{i}pio receber incentivos financeiros ($implantados\_esf$), dados do Departamento de Aten\c{c}\~{a}o B\'{a}sica (DAB). 

Devido a existência de uma quantidade consider\'{a}vel de zeros em algumas dessas vari\'{a}veis, foi utilizada a transformação proposta por \citeonline{bibJohnson1949} na aplicação do logaritmo natural:

\begin{equation}
 f(y,\theta) = \text{sinh}^{-1}(\theta y) / \theta = \frac{\log{( \theta y + (\theta^2 y^2 + 1)^{1/2} )}}{\theta}
 \label{eq_log_transform}
\end{equation}
\noindent com $\theta = 0.5$.

Por fim, a Tabela \ref{tabEstatisticas} mostra algumas estat\'{i}sticas descritivas para variáveis de interesse que atendem aos resultados parciais da pesquisa. Cabe ressaltar que as variáveis reais foram deflacionadas pelo índice IGP-DI da Fundação Getulio Vargas, \citeonline{bibFGV2016}, com ano base 2010.

\begin{table}[htbp]\centering \caption{Estat\'{i}sticas das vari\'{a}veis \label{tabEstatisticas}}
\begin{tabular}{l c c c c c}\hline\hline
\multicolumn{1}{c}{\textbf{Vari\'{a}vel}} & \textbf{M\'{e}dia} & \textbf{Desv. Pad.}& \textbf{Min.} &  \textbf{M\'{a}x.} & \textbf{N} \\ \hline
txprenatal & 64.088 & 19.727 & 0.511 & 100 & 16686\\
%pib2realpc & 13.928 & 24.49 & -89.497 & 2281.565 & 16686\\
pibrealpc & 13.645 & 16.626 & 1.256 & 725.275 & 16686\\
distorcao\_ef & 23.287 & 11.655 & 0.4 & 66.099 & 16682\\
abandono\_em & 9.377 & 6.72 & 0 & 100 & 16650\\
aguatratada & 9030.541 & 41277.313 & 0 & 2205269 & 16686\\
coletalixo & 4917.458 & 23771.949 & 0 & 1296401 & 16686\\
despesaspsreal & 8192017.579 & 66354439.674 & 0 & 4382966272 & 16686\\
equipamentos & 214.854 & 2143.256 & 0 & 137830 & 16686\\
implantados\_esf & 5.762 & 15.619 & 0 & 757 & 16686\\
\hline
\end{tabular}
\begin{tablenotes}
\item \tiny{Fonte: Elaboração dos autores.}
\end{tablenotes}
\end{table}


% \begin{table}[!htb]
% \caption{Descri\c{c}\~{a}o das vari\'{a}veis da base de dados \label{tabDescricaoBD}}
% \centering
% {\footnotesize
% \begin{tabular}{cp{4cm}p{9cm}}
% \toprule\toprule
% \textbf{Vari\'{a}vel} & \textbf{Descri\c{c}\~{a}o} & \textbf{Fonte} \\
% \midrule
% \midrule
% $y_2$ & Descri\c{c}\~{a}o da vari\'{a}vel $y_2$ & \multirow{2}{*}{ \begin{minipage}{9cm} Datasus - Sistema de Informa\c{c}\~{a}o de Aten\c{c}\~{a}o B\'{a}sica. \url{http://tabnet.datasus.gov.br/cgi/deftohtm.exe?siab/cnv/SIABSBR.DEF} \end{minipage} } \\[1mm]
% \cmidrule{1-2}
% $y_1$ & Descri\c{c}\~{a}o da vari\'{a}vel $y_1$ & \\ 
% \midrule
% $txmortalidade$ & Taxa de mortalidade & \begin{equation} txmortalidade = \frac{obitoinfati}{nvpn} \times 100 \label{eqTxmortalidade} \end{equation} \\
% \midrule
% $poptcu$ & Popula\c{c}\~{a}o usada para fins de c\'{a}lculo do TCU & \multirow{1}{*}{ \begin{minipage}{9cm} Datasus - Popula\c{c}\~{a}o Residente, Estimativas para o TCU. \url{http://tabnet.datasus.gov.br/cgi/deftohtm.exe?ibge/cnv/poptbr.def} \end{minipage} } \\[8mm]
% \bottomrule
% \end{tabular}
% }
% \end{table}

%\subsubsection{Problemas com os dados}
\subsubsection{Limitações dos dados}

Deve-se destacar algumas limita\c{c}\~{o}es dos dados advindos do Siab. Trata-se de dados gerados a partir do trabalho das eSF e Agentes Comunit\'{a}rios de Sa\'{u}de (ACS), abrangendo a situa\c{c}\~{a}o de saneamento,  moradia e sa\'{u}de das fam\'{i}lias. Essas informa\c{c}\~{o}es excluem os munic\'{i}pios que n\~{a}o informaram todos os meses do per\'{i}odo selecionado, raz\~{a}o pela qual se poder\'{a} ter indicadores diferentes no cruzamento das vari\'{a}veis aqui disponibilizadas. Aplica-se tamb\'{e}m uma rotina para a cr\'{i}tica dos dados. Esta rotina se baseia na defini\c{c}\~{a}o de crit\'{e}rios, a partir dos quais se define pela inclus\~{a}o ou exclus\~{a}o do munic\'{i}pio na base de dados para an\'{a}lise - ``base limpa''. O \'{o}rg\~{a}o respons\'{a}vel definiu crit\'{e}rios de verifica\c{c}\~{a}o de erros e inconsist\^{e}ncias, tanto para a base de dados de cadastro quanto para a base de dados de situa\c{c}\~{a}o de sa\'{u}de. Por\'{e}m, ainda n\~{a}o foram definidos crit\'{e}rios para limpeza da base de dados de produ\c{c}\~{a}o.

Ap\'{o}s aplica\c{c}\~{a}o das rotinas de limpeza, obt\^{e}m-se duas ``bases limpas'': a base de cadastro e a base de situa\c{c}\~{a}o de sa\'{u}de. A ``base limpa'' de cadastro exclui os munic\'{i}pios com erros ou inconsist\^{e}ncias relacionados a qualquer um dos crit\'{e}rios considerados. A ``base limpa'' de situa\c{c}\~{a}o de sa\'{u}de inclui ou exclui o munic\'{i}pio com rela\c{c}\~{a}o a cada um dos indicadores analisados. Vale ressaltar que, como as rotinas s\~{a}o independentes, a exclus\~{a}o de um munic\'{i}pio numa das bases n\~{a}o implica na sua exclus\~{a}o da outra base.

Existe ainda o problema de dados sujos. A Tabela \ref{tabBdSujo_obitosdoenca} mostra alguns exemplos de registros incomuns do banco de dados de \'{o}bitos por motivos de doen\c{c}a. O ano de 2013 apresenta dados muito superiores aos dos anos anteriores. Nos registros em que se observou esse tipo de ocorrência, optou-se por desconsiderar o valor do ano e recalcul\'{a}-lo por meio de uma regress\~{a}o semelhante à equação \ref{eq_estimativa_pib_municipal_1}.

\begin{table}[htbp]\centering \caption{Alguns registros sujos do banco de dados de \'{o}bitos por motivos de doen\c{c}a. O ano de 2013 apresenta dados muito superiores aos dos anos anteriores. \label{tabBdSujo_obitosdoenca}}
\begin{tabular}{c l c c c c c c c}\hline\hline
\multicolumn{2}{c}{\textbf{Munic\'{i}pio}} & ~ & \multicolumn{6}{c}{\textbf{obitosdoenca}} \\ \cline{1-2} \cline{4-9}
\textbf{C\'{o}digo} & \textbf{Nome} & ~ & \textbf{2008} & \textbf{2009} & \textbf{2010} & \textbf{2011} & \textbf{2012} & \textbf{2013} \\
\hline
290010 & Aba\'{i}ra & ~ & 4 & 4 & 7 & 6 & 5 & 4029 \\
351515 & Engenheiro Coelho & ~ & 11 & 19 & 16 & 18 & 21 & 6541 \\
430160 & Bag\'{e} & ~ & 5 & 7 & 6 & 6 & 6 & 283 \\
430210 & Bento Gon\c{c}alves & ~ & 1 & 1 & 2 & 6 & 2 & 148 \\
430350 & Camaqu\~{a} & ~ & 12 & 8 & 6 & 5 & 6 & 136 \\
430510 & Caxias do Sul & ~ & 18 & 22 & 19 & 22 & 11 & 647 \\
430920 & Gravata\'{i} & ~ & 4 & 8 & 9 & 6 & 8 & 534 \\
431440 & Pelotas & ~ & 6 & 6 & 2 & 5 & 7 & 692 \\
432130 & Taquari & ~ & 65 & 71 & 67 & 59 & 53 & 12784 \\
\hline
\end{tabular}
\begin{tablenotes}
\item \tiny{Fonte: Elaboração dos autores.}
\end{tablenotes}
\end{table}

Al\'{e}m disso, foram detectados alguns poucos municípios com valores do PIB com variação incomum, conforme a Tabela \ref{tabBdSujo_pib}. Não sendo porém possível atribuir essas variações a alguma sazonalidade, decidiu-se pela exclusão dos mesmos. Ademais, a despesa com recurso pr\'{o}prio do munic\'{i}pio Jucuru\c{c}u, c\'{o}digo 2918456, consta como negativa para o ano de 2012.  Neste caso, considerou-se o seu valor absoluto, pois a magnitude est\'{a} adequada no contexto dos outros valores anuais.

\begin{table}[htbp]\centering \caption{Alguns registros sujos do banco de dados de PIB municipal. Existe uma varia\c{c}\~{a}o irreal entre os anos. \label{tabBdSujo_pib}}
\begin{tabular}{c l c c c c c c}\hline\hline
\multicolumn{2}{c}{\textbf{Munic\'{i}pio}} & ~ & \multicolumn{5}{c}{\textbf{pib}} \\ \cline{1-2} \cline{4-8}
\textbf{C\'{o}digo} & \textbf{Nome} & ~ & \textbf{2010} & \textbf{2011} & \textbf{2012} & \textbf{2013} & \textbf{2014} \\
\hline
2404507 & Guamar\'{e} & ~ & 631523.5 & 245576.2 & 11071.7 & 52826 & 240450 \\
2929206 & S\~{a}o Francisco do Conde & ~ & 5659444 & 2882176 & 1842066 & 1566375 & 292920 \\
4304358 & Candiota & ~ & 612573.4 & 344585.6 & 129429.5 & 137829.4 & 430435 \\
\hline
\end{tabular}
\begin{tablenotes}
\item \tiny{Fonte: Elaboração dos autores.}
\end{tablenotes}
\end{table}


\subsection{Estrat\'{e}gias emp\'{i}ricas}\label{secEstrategiaEmp}

O proposta de identificar um melhor indicador para saúde humana está além dos resultados parciais apresentados aqui. Por isso, utilizou-se resultados preliminares para indicar objetivamente que a vari\'{a}vel dependente de interesse \'{e} a propor\c{c}\~{a}o de nascidos vivos de m\~{a}es com sete ou mais consultas de pr\'{e}-natal ($txprenatal$). Definida conforme a Tabela \ref{tabDVarDependente}, ela equaciona a quantidade de nascimentos com vida de gestantes que obtiveram mais do que 7 consultas antes do nascimento do beb\^{e} ($nv7oumaisconsult$) e a quantidade total de nascimentos com vida, independente da quantidade de consultas que antecederam o nascimento do beb\^{e} ($nvpn$). A hipótese subjacente é de que essa variável é representativa do estado de saúde da população por abstrair sobre impacto dos cuidados médicos sobre o estado de saúde da mãe e da criança.

Assim, o efeito do Programa Mais M\'{e}dicos sobre o estado de saúde da população ser\'{a} obtido atrav\'{e}s da estima\c{c}\~{a}o de diferen\c{c}as em diferen\c{c}as dos seguintes modelos:
\begin{enumerate}
 \item Uma estima\c{c}\~{a}o de diferen\c{c}as em diferen\c{c}as simples:
  \begin{equation}
  txprenatal_{it} = \beta_0 + \beta_1 \, (PMM)_{i} + \beta_2 \, (d2014)_{t} + \beta_3 \, (PMM \cdot d2014)_{it} + \varepsilon_{it}
  \label{eq_mod_txprenatal}
  \end{equation}
  
  \noindent onde $PMM$ \'{e} uma vari\'{a}vel \textit{dummy} que indica se o Munic\'{i}pio participa do Programa Mais M\'{e}dicos, e
  $d2014$ \'{e} uma vari\'{a}vel \textit{dummy} igual a 1 se o ano for 2014.

 \item Uma estima\c{c}\~{a}o de diferen\c{c}as em diferen\c{c}as com covariadas:
  \begin{description}
   \item[a)] socioeconômicas:
    \begin{equation}
    \begin{array}{rcl}
     txprenatal_{it} & = & \beta_0 + \beta_1 \, (PMM)_{i} + \beta_2 \, (d2014)_{t} + \beta_3 \, (PMM \cdot d2014)_{it} \\
     ~               & ~ & + \beta_4 \, (pibrealpc)_{it} + \beta_5 \, (distorcao\_ef)_{it} \\
     ~               & ~ & + \beta_6 \, (abandono\_em)_{it} + \varepsilon_{it}
    \end{array}
    \label{eq_mod_txprenatal_cov1}
    \end{equation}
    \noindent onde 
    $pibrealpc$ \'{e} o logar\'{i}tmo do PIB real per capita, com o ano de 2014 estimado pela eq. (\ref{eq_estimativa_pib_municipal_1}),
    $distorcao\_ef$ \'{e} a taxa de distor\c{c}\~{a}o idade-s\'{e}rie total no ensino fundamental, e
    $abandono\_em$ \'{e} a taxa de abandono total no ensino m\'{e}dio.

    \item[b)] ambientais:
    \begin{equation}
    \begin{array}{rcl}
     txprenatal_{it} & = & \beta_0 + \beta_1 \, (PMM)_{i} + \beta_2 \, (d2014)_{t} + \beta_3 \, (PMM \cdot d2014)_{it} \\
     ~               & ~ & + \beta_4 \, (\log{(aguatratada)})_{it} + \beta_5 \, (\log{(coletalixo)})_{it} + \varepsilon_{it}
    \end{array}
    \label{eq_mod_txprenatal_cov2}
    \end{equation}
    \noindent onde $\log{(aguatratada)}$ \'{e} o logar\'{i}tmo de $aguaclorada + aguafervida + aguafiltrada + aguapublica$,
    $aguaclorada$, $aguafervida$ e $aguafiltrada$ s\~{a}o os domic\'{i}lios com \'{a}gua clorara, fervida e filtrada, respectivamente,
    $aguapublica$ s\~{a}o os domic\'{i}lios servidos de \'{a}gua proveniente de uma rede geral,
    $\log{(coletalixo)}$ \'{e} o logar\'{i}tmo de dos domic\'{i}lios com lixo coletado.
    
    \item[c)] s\'{o}cio-econ\^{o}micas, ambientais e relacionadas a servi\c{c}os m\'{e}dicos:
    \begin{equation}
    \begin{array}{rcl}
     txprenatal_{it} & = & \beta_0 + \beta_1 \, (PMM)_{i} + \beta_2 \, (d2014)_{t} + \beta_3 \, (PMM \cdot d2014)_{it} \\
     ~               & ~ & + \beta_4 \, (pibrealpc)_{it} + \beta_5 \, (distorcao\_ef)_{it} + \beta_6 \, (abandono\_em)_{it} \\
     ~               & ~ & + \beta_7 \, (\log{(aguatratada)})_{it} + \beta_8 \, (\log{(coletalixo)})_{it} \\
     ~               & ~ & + \beta_9 \, (\log{(despesaspsreal)})_{it} + \beta_{10} \, (equipamentos)_{it} \\
     ~               & ~ &  + \beta_{11} \, (\log{(implantados\_esf)})_{it} + \varepsilon_{it}\\
    \end{array}
    \label{eq_mod_txprenatal_cov3}
    \end{equation}
    \noindent onde 
    $\log{(despesaspsreal)}$ \'{e} o logar\'{i}tmo das despesas com recursos pr\'{o}prios deflacionadas,
    $equipamentos$ s\~{a}o os equipamentos existentes em uso pelo SUS, e \\
    $\log{(implantados\_esf)}$ \'{e} o n\'{u}mero de eSF aptas para o munic\'{i}pio receber incentivos financeiros.

 \item[d)] Uma estima\c{c}\~{a}o de diferen\c{c}as em diferen\c{c}as, equação (\ref{eq_mod_txprenatal}),
  segmentada pela renda per capita:
   \begin{enumerate}
    \item $N$ munic\'{i}pios com a menor renda per capita;
    \item o restante dos munic\'{i}pios.
   \end{enumerate}
 
  \end{description}
\end{enumerate}

Para cada modelo proposto, ser\~{a}o feitas tr\^{e}s regress\~{o}es: o teste de tend\^{e}ncia temporal proposto na se\c{c}\~{a}o \ref{secTendenciaNoTempo} para verificar se a hip\'{o}tese de tend\^{e}ncia no tempo comum \'{e} atendida, uma regress\~{a}o por OLS para fins comparativos, e uma regress\~{a}o por Efeitos Fixos de forma a garantir a elimina\c{c}\~{a}o de efeitos individuais n\~{a}o-observ\'{a}veis.

%------------------------------------------------

\section{Resultados \label{secResultados}}

Os resultados apresentados a seguir estão de acordo com as estratégias empíricas da seção \ref{secEstrategiaEmp}. 
%O programa utilizado para gerar os resultados desta se\c{c}\~{a}o foi o Stata, vers\~{a}o 12, para Linux. Temos uma versão original?
% O Ap\^{e}ndice \ref{secApendice_saidas} mostra as sa\'{i}das originais que ser\~{a}o mostrados a seguir. 

\subsection{Diferen\c{c}as em diferen\c{c}as simples}

Inicialmente, obteve-se o teste de tendência comum da variável dependente "taxa pr\'{e}-natal" tanto em termos gráficos quanto com base na estimação proposta na seção \ref{secTendenciaNoTempo}. Visualmente, a Figura \ref{fig_txprenatal_graf_tend} mostra que a evolu\c{c}\~{a}o temporal dos grupos de tratamento e de controle seguem a mesma tend\^{e}ncia entre 2010 e 2012, mas o mesmo n\~{a}o ocorre entre 2012 e 2014. O teste de tend\^{e}ncia corrobora com este resultado. Como mostrado na coluna 1 da Tabela \ref{tab_txprenatal}, a hipótese $H_0: \text{PMM\_d2012} = 0$ é aceita, ao passo que rejeita-se $H_0: \text{PMM\_d2014} = 0$. Isso implica que, de fato, a hip\'{o}tese de tend\^{e}ncia no tempo comum \'{e} atendida. Logo, os grupos de tratados e controle tinham a mesma tend\^{e}ncia antes da intervenção, evidenciando a existência de efeito real do programa sobre o estado de saúde da população.

\begin{figure}[!htb]
\begin{center}
 \includegraphics[width=0.8\textwidth]{figuras/fig_txprenatal_graf_tend.png}
\caption{\label{fig_txprenatal_graf_tend} Evolu\c{c}\~{a}o da taxa de pr\'{e}-natal, eq. (\ref{eq_txprenatal}), ao longo do tempo.
(Azul) $E[txprenatal | ano, PMM = 0]$, taxa de pr\'{e}-natal para os controles;
e (Vermelho) $E[txprenatal | ano, PMM = 1]$, taxa de pr\'{e}-natal para os tratados.}
\end{center}
\end{figure}

\begin{table}[htbp]\centering
\def\sym#1{\ifmmode^{#1}\else\(^{#1}\)\fi}
\caption{Diferen\c{c}as em diferen\c{c}as da vari\'{a}vel \textit{txprenatal}, equação (\ref{eq_mod_txprenatal}). (1) Teste de tend\^{e}ncia; (2) OLS; (3) FE.\label{tab_txprenatal}}
\begin{tabular*}{0.9\hsize}{@{\hskip\tabcolsep\extracolsep\fill}l*{3}{D{.}{.}{-1}}}
\toprule
            &\multicolumn{1}{c}{(1)}&\multicolumn{1}{c}{(2)}&\multicolumn{1}{c}{(3)}\\
            &\multicolumn{1}{c}{txprenatal}&\multicolumn{1}{c}{txprenatal}&\multicolumn{1}{c}{txprenatal}\\
            &\multicolumn{1}{c}{$\beta$ / Erro Padr\~{a}o}   &\multicolumn{1}{c}{$\beta$ / Erro Padr\~{a}o}   &\multicolumn{1}{c}{$\beta$ / Erro Padr\~{a}o}   \\
\midrule
PMM         &      -7.781***&      -6.698***&               \\
            &     (0.554)   &     (0.507)   &               \\
d2012       &       2.485***&               &               \\
            &     (0.645)   &               &               \\
d2014       &       4.850***&       2.365***&       2.365***\\
            &     (0.645)   &     (0.591)   &     (0.243)   \\
PMM\_d2012   &       1.083   &               &               \\
            &     (0.783)   &               &               \\
PMM\_d2014   &       2.843***&       1.760** &       1.760***\\
            &     (0.783)   &     (0.717)   &     (0.295)   \\
Constante   &      66.032***&      68.517***&      63.975***\\
            &     (0.456)   &     (0.418)   &     (0.097)   \\
\midrule
F           &\multicolumn{1}{c}{153.4}   &\multicolumn{1}{c}{127.3}   &\multicolumn{1}{c}{351.7}   \\
\(R^{2}\)   &\multicolumn{1}{c}{0.044}   &\multicolumn{1}{c}{0.033}   &\multicolumn{1}{c}{0.112}   \\
N           &\multicolumn{1}{c}{16686}   &\multicolumn{1}{c}{11124}   &\multicolumn{1}{c}{11124}   \\
\bottomrule
\multicolumn{4}{l}{\footnotesize * $p<0.10$, ** $p<0.05$, *** $p<0.01$}\\
\end{tabular*}
\begin{tablenotes}
\item \tiny{Fonte: Elaboração dos autores.}
\end{tablenotes}
\end{table}

As colunas 2 e 3 da Tabela \ref{tab_txprenatal} mostram as regress\~{o}es por OLS e por efeitos fixos que estimam o efeito do tratamento
por diferen\c{c}as em diferen\c{c}as. O efeito do tratamento \'{e} positivo, igual a $1.755$, e estatisticamente significante ao n\'{i}vel de $1\%$.
O sinal \'{e} condizente com o previsto pela teoria, conforme mostrado na Tabela \ref{tabDVarDependente}. Observe que a hip\'{o}tese de tend\^{e}ncia no tempo comum \'{e} atendida para todos os modelos, pois $\text{PMM\_d2012}$ n\~{a}o \'{e} estatisticamente significante. Logo, \'{e} poss\'{i}vel realizar a estimativa por diferen\c{c}as em diferen\c{c}as.


\subsection{Diferen\c{c}as em diferen\c{c}as usando covariadas}

As estimativas por diferen\c{c}as em diferen\c{c}as usando covariadas
socioeconômicas, ambientais e relacionadas a servi\c{c}os m\'{e}dicos s\~{a}o mostradas nas
Tabelas \ref{eq_mod_txprenatal_cov1}, \ref{eq_mod_txprenatal_cov2} e \ref{eq_mod_txprenatal_cov3},
respectivamente.
As colunas 2 e 3 das Tabelas \ref{tab_txprenatal_cov1}, \ref{tab_txprenatal_cov2} e \ref{tab_txprenatal_cov3}
mostram as regress\~{o}es por OLS e por efeitos fixos que estimam o efeito do tratamento
por diferen\c{c}as em diferen\c{c}as usando as equa\c{c}\~{o}es (\ref{eq_mod_txprenatal_cov1}),
(\ref{eq_mod_txprenatal_cov2}) e (\ref{eq_mod_txprenatal_cov3}), respectivamente.
Os efeitos do tratamento, iguais a $1.722$, $1.831$ e $1.725$, s\~{a}o
estatisticamente significantes ao n\'{i}vel de $1\%$.
Os valores s\~{a}o bem pr\'{o}ximos do encontrado sem o uso de covariadas,
sugerindo que o efeito \'{e} consistente.

\begin{table}[htb]\centering
\def\sym#1{\ifmmode^{#1}\else\(^{#1}\)\fi}
\caption{Diferen\c{c}as em diferen\c{c}as da vari\'{a}vel \textit{txprenatal} usando covariadas s\'{o}cio-econ\^{o}micas, equação  (\ref{eq_mod_txprenatal_cov1}). (1) Teste de tend\^{e}ncia; (2) OLS; (3) FE.\label{tab_txprenatal_cov1}}
\begin{tabular*}{0.9\hsize}{@{\hskip\tabcolsep\extracolsep\fill}l*{3}{D{.}{.}{-1}}}
\toprule
            &\multicolumn{1}{c}{(1)}&\multicolumn{1}{c}{(2)}&\multicolumn{1}{c}{(3)}\\
            &\multicolumn{1}{c}{txprenatal}&\multicolumn{1}{c}{txprenatal}&\multicolumn{1}{c}{txprenatal}\\
            &\multicolumn{1}{c}{$\beta$ / Erro Padr\~{a}o}   &\multicolumn{1}{c}{$\beta$ / Erro Padr\~{a}o}   &\multicolumn{1}{c}{$\beta$ / Erro Padr\~{a}o}   \\
\midrule
PMM         &      -4.203***&      -3.581***&               \\
            &     (0.437)   &     (0.410)   &               \\
d2012       &       0.214   &               &               \\
            &     (0.508)   &               &               \\
d2014       &      -0.420   &      -0.373   &       1.850***\\
            &     (0.510)   &     (0.477)   &     (0.272)   \\
PMM\_d2012   &       0.970   &               &               \\
            &     (0.616)   &               &               \\
PMM\_d2014   &       2.683***&       1.721***&       1.759***\\
            &     (0.616)   &     (0.577)   &     (0.295)   \\
pibrealpc   &       0.082***&       0.056***&      -0.056*  \\
            &     (0.007)   &     (0.008)   &     (0.032)   \\
distorcao\_ef&      -0.917***&      -0.837***&      -0.176***\\
            &     (0.012)   &     (0.014)   &     (0.044)   \\
abandono\_em &      -0.321***&      -0.304***&      -0.063***\\
            &     (0.020)   &     (0.023)   &     (0.023)   \\
Constante   &      89.445***&      88.188***&      69.489***\\
            &     (0.468)   &     (0.471)   &     (1.132)   \\
\midrule
F           &\multicolumn{1}{c}{1451.6}   &\multicolumn{1}{c}{1118.8}   &\multicolumn{1}{c}{146.6}   \\
\(R^{2}\)   &\multicolumn{1}{c}{0.411}   &\multicolumn{1}{c}{0.377}   &\multicolumn{1}{c}{0.117}   \\
N           &\multicolumn{1}{c}{16646}   &\multicolumn{1}{c}{11098}   &\multicolumn{1}{c}{11098}   \\
\bottomrule
\multicolumn{4}{l}{\footnotesize * $p<0.10$, ** $p<0.05$, *** $p<0.01$}\\
\end{tabular*}
\begin{tablenotes}
\item \tiny{Fonte: Elaboração dos autores.}
\end{tablenotes}
\end{table}

\begin{figure}[!htb]
\begin{center}
 \includegraphics[width=0.8\textwidth]{figuras/fig_qreg_txprenatal_cov1.png}
\caption{\label{fig_qreg_txprenatal_cov1} Regress\~{a}o quant\'{i}lica do modelo com covariadas s\'{o}cio-econ\^{o}micas, eq. (\ref{eq_mod_txprenatal_cov1}).}
\end{center}
\end{figure}


\begin{table}[htb]\centering
\def\sym#1{\ifmmode^{#1}\else\(^{#1}\)\fi}
\caption{Diferen\c{c}as em diferen\c{c}as da vari\'{a}vel \textit{txprenatal} usando covariadas ambientais, equação (\ref{eq_mod_txprenatal_cov2}). (1) Teste de tend\^{e}ncia; (2) OLS; (3) FE.\label{tab_txprenatal_cov2}}
\begin{tabular*}{0.9\hsize}{@{\hskip\tabcolsep\extracolsep\fill}l*{3}{D{.}{.}{-1}}}
\toprule
            &\multicolumn{1}{c}{(1)}&\multicolumn{1}{c}{(2)}&\multicolumn{1}{c}{(3)}\\
            &\multicolumn{1}{c}{txprenatal}&\multicolumn{1}{c}{txprenatal}&\multicolumn{1}{c}{txprenatal}\\
            &\multicolumn{1}{c}{$\beta$ / Erro Padr\~{a}o}   &\multicolumn{1}{c}{$\beta$ / Erro Padr\~{a}o}   &\multicolumn{1}{c}{$\beta$ / Erro Padr\~{a}o}   \\
\midrule
PMM         &      -6.773***&      -5.623***&               \\
            &     (0.524)   &     (0.486)   &               \\
d2012       &       2.366***&               &               \\
            &     (0.608)   &               &               \\
d2014       &       4.774***&       2.433***&       2.277***\\
            &     (0.619)   &     (0.573)   &     (0.260)   \\
PMM\_d2012   &       1.187   &               &               \\
            &     (0.738)   &               &               \\
PMM\_d2014   &       2.958***&       1.782** &       1.831***\\
            &     (0.751)   &     (0.695)   &     (0.314)   \\
log\_aguatratada&     -15.414***&     -13.676***&      -2.378***\\
             &     (0.292)   &     (0.343)   &     (0.883)   \\
log\_coletalixo&      13.930***&      12.168***&       2.352***\\
            &     (0.272)   &     (0.323)   &     (0.813)   \\
Constante   &      87.440***&      88.775***&      65.752***\\
            &     (1.074)   &     (1.192)   &     (2.816)   \\
\midrule
F           &\multicolumn{1}{c}{535.1}   &\multicolumn{1}{c}{402.6}   &\multicolumn{1}{c}{158.4}   \\
\(R^{2}\)   &\multicolumn{1}{c}{0.193}   &\multicolumn{1}{c}{0.163}   &\multicolumn{1}{c}{0.114}   \\
N           &\multicolumn{1}{c}{15716}   &\multicolumn{1}{c}{10363}   &\multicolumn{1}{c}{10363}   \\
\bottomrule
\multicolumn{4}{l}{\footnotesize * $p<0.10$, ** $p<0.05$, *** $p<0.01$}\\
\end{tabular*}
\begin{tablenotes}
\item \tiny{Fonte: Elaboração dos autores.}
\end{tablenotes}
\end{table}

% \begin{figure}[!htb]
% \begin{center}
%  \includegraphics[width=0.8\textwidth]{figuras/fig_qreg_txprenatal_cov2.png}
% \caption{\label{fig_qreg_txprenatal_cov2} Regress\~{a}o quant\'{i}lica do modelo com covariadas ambientais, eq. (\ref{eq_mod_txprenatal_cov2}).}
% \end{center}
% \end{figure}


\begin{table}[!htb]\centering
\def\sym#1{\ifmmode^{#1}\else\(^{#1}\)\fi}
\caption{Diferen\c{c}as em diferen\c{c}as da vari\'{a}vel \textit{txprenatal} usando covariadas s\'{o}cio-econ\^{o}micas, ambientais e relacionadas a servi\c{c}os m\'{e}dicos, equação (\ref{eq_mod_txprenatal_cov3}). (1) Teste de tend\^{e}ncia; (2) OLS; (3) FE.\label{tab_txprenatal_cov3}}
\begin{tabular*}{0.95\hsize}{@{\hskip\tabcolsep\extracolsep\fill}l*{3}{D{.}{.}{-1}}}
\toprule
            &\multicolumn{1}{c}{(1)}&\multicolumn{1}{c}{(2)}&\multicolumn{1}{c}{(3)}\\
            &\multicolumn{1}{c}{txprenatal}&\multicolumn{1}{c}{txprenatal}&\multicolumn{1}{c}{txprenatal}\\
            &\multicolumn{1}{c}{$\beta$ / Erro Padr\~{a}o}   &\multicolumn{1}{c}{$\beta$ / Erro Padr\~{a}o}   &\multicolumn{1}{c}{$\beta$ / Erro Padr\~{a}o}   \\
\midrule
PMM         &      -6.551***&      -5.738***&               \\
            &     (0.485)   &     (0.455)   &               \\
d2012       &       1.480***&               &               \\
            &     (0.563)   &               &               \\
d2014       &       1.514***&       0.356   &       2.211***\\
            &     (0.566)   &     (0.531)   &     (0.251)   \\
PMM\_d2012   &       1.051   &               &               \\
            &     (0.683)   &               &               \\
PMM\_d2014   &       2.845***&       1.797***&       1.731***\\
            &     (0.683)   &     (0.641)   &     (0.296)   \\
pibrealpc   &       0.115***&       0.087***&      -0.054*  \\
            &     (0.008)   &     (0.009)   &     (0.032)   \\
abandono\_em &      -0.826***&      -0.730***&      -0.060***\\
            &     (0.020)   &     (0.023)   &     (0.023)   \\
log\_aguatratada&      -4.988***&      -4.176***&      -0.462*  \\
            &     (0.127)   &     (0.150)   &     (0.252)   \\
log\_coletalixo&       5.038***&       4.282***&       0.493*  \\
            &     (0.134)   &     (0.159)   &     (0.273)   \\
log\_despesaspsreal&       0.375***&       0.382***&       0.092***\\
            &     (0.043)   &     (0.045)   &     (0.033)   \\
equipamentos&       0.000***&       0.000***&      -0.001** \\
            &     (0.000)   &     (0.000)   &     (0.000)   \\
implantados\_esf&      -0.106***&      -0.106***&       0.107** \\
            &     (0.010)   &     (0.011)   &     (0.044)   \\
Constante   &      68.417***&      67.366***&      62.475***\\
            &     (1.340)   &     (1.401)   &     (1.221)   \\
\midrule
F           &\multicolumn{1}{c}{527.0}   &\multicolumn{1}{c}{334.6}   &\multicolumn{1}{c}{81.94}   \\
\(R^{2}\)   &\multicolumn{1}{c}{0.275}   &\multicolumn{1}{c}{0.232}   &\multicolumn{1}{c}{0.117}   \\
N           &\multicolumn{1}{c}{16650}   &\multicolumn{1}{c}{11102}   &\multicolumn{1}{c}{11102}   \\
\bottomrule
\multicolumn{4}{l}{* $p<0.10$, ** $p<0.05$, *** $p<0.01$}\\
\end{tabular*}
\begin{tablenotes}
\item \tiny{Fonte: Elaboração dos autores.}
\end{tablenotes}
\end{table}


% \begin{figure}[!htb]
% \begin{center}
%  \includegraphics[width=0.8\textwidth]{figuras/fig_qreg_txprenatal_cov3.png}
% \caption{\label{fig_qreg_txprenatal_cov3} Regress\~{a}o quant\'{i}lica do modelo com covariadas s\'{o}cio-econ\^{o}micas, ambientais e relacionadas a servi\c{c}os m\'{e}dicos, eq. (\ref{eq_mod_txprenatal_cov3}).}
% \end{center}
% \end{figure}


O modelo na forma geral da seção \ref{secModelosSaude} fornece uma estrutura te\'{o}rica que permite analisar os sinais dos coeficientes das covariadas.
Com rela\c{c}\~{a}o ao resultado para o modelo com as covariadas socioeconômicas (Tabela \ref{tab_txprenatal_cov1}),
a teoria prediz que $\beta_{pibrealpc}$ deve ser positivo, mas
a regress\~{a}o por OLS retorna $\beta_{pibrealpc} = 0.056 > 0$ (estatisticamente significante a 1\%)
e a regress\~{a}o por FE retorna $\beta_{pibrealpc} = -0.056 < 0$ (estatisticamente significante a 10\%),
o que \'{e} conflitante. Considerando que saúde humana é um bem de luxo, um efeito renda nulo é razoável para faixas de rendas mais baixas enquanto que deve haver efeito positivo para faixas mais elevadas. Além de justificar a análise em nível de renda, este resultado aponta para a necessidade de se averiguar os fatos empíricos que explicam os resultados aqui obtidos.

Os coeficientes de $distorcao\_ef$ e $abandono\_em$ s\~{a}o ambos negativos e estatisticamente significantes ao n\'{i}vel de 1\%. Note que as \textit{proxies} aqui utilizadas para educação formal estão de acordo com o efeito esperado pela teoria, municípios com menores índices de evasão escolar e distorção na idade-série do aluno, tende a impactar positivamente a taxa pré-natal. Os coeficientes da regress\~{a}o quant\'{i}lica da equação (\ref{eq_mod_txprenatal_cov1}), dispostos na Figura \ref{fig_qreg_txprenatal_cov1}, mostram que esses sinais se mantiveram ao longo de todos os percentis entre 10\% e 90\%, mesmo considerando os intervalos de confian\c{c}a.

Com rela\c{c}\~{a}o ao resultado para o modelo com as covariadas ambientais (Tabela \ref{tab_txprenatal_cov2}), apesar dos dados estarem ajustados estatisticamente, apenas o coeficiente de $\log{(coletalixo)}$ \'{e} positivo e está de acordo com a expectativa a priori. O resultado negativo de $\log{(aguatratada)}$ \'{e} inesperado e cabe pesquisas adicionais que explorem esta relação. A ausência de mecanismo de transmissão entre tais condições ambientais e a taxa pré-natal pode justificar tanto o resultado inesperado quanto a pouca relevância dessas vari\'{a}veis para explicar a vari\'{a}vel de interesse.

Os resultados para o modelo com as covariadas socioeconômicas, ambientais e relacionados a serviços médicos estão na Tabela \ref{tab_txprenatal_cov3}. Os coeficientes de $\log{(despesaspsreal)}$ e de ($implantados\_esf$) s\~{a}o positivos e estão de acordo com o esperado. O mesmo não se verifica para a variável $equipamentos$, que \'{e} negativo. 


\subsection{Diferen\c{c}as em diferen\c{c}as segmentado por renda per capita}

Objetivando compreender o efeito renda por segmentos, estimou-se o modelo em diferenças em diferenças básico por faixas de renda per capita. A Figura \ref{fig_txprenatal_efeito_segmentado} mostra o efeito do programa sobre $txprenatal$ com base no coeficiente de $PMM \cdot d2014$ da equação (\ref{eq_mod_txprenatal}) segmentado para diferentes valores de $N$.
As regress\~{o}es foram feitas em dois segmentos, a primeira sobre os $N$ munic\'{i}pios com menor renda per capita
e a segunda sobre os munc\'{i}pios restantes. As m\'{e}dias para cada um destes segmentos são mostradas na Figura \ref{fig_txprenatal_media_segmentado}.

\begin{figure}[!htb]
\begin{center}
\begin{overpic}[width=14cm]{figuras/fig_txprenatal_efeito_segmentado.png}
\put(0,40){\rotatebox{90}{Efeito}}
\put(50,2){N}
\end{overpic}
%\caption{\label{fig_txprenatal_efeito_segmentado} Efeito do Programa Mais M\'{e}dicos sobre $txprenatal$, coeficiente de $PMM \cdot d2014$ na eq. (\ref{eq_mod_txprenatal}),
%segmentado pela renda per capita. A linha vermelha corresponde ao efeito com o segmento dos $N$ munic\'{i}pios com menor renda per capita,
%e a linha azul corresponde o efeito do segmento composto pelo restante dos munic\'{i}pios.
%Pontos azul e vermelho claros s\~{a}o estat\'{i}sticamente significantes ao n\'{i}vel de 10\%,
%enquanto os pontos escuros n\~{a}o s\~{a}o.}
\caption{\label{fig_txprenatal_efeito_segmentado} Efeito do Programa Mais M\'{e}dicos sobre $txprenatal$ segmentado por renda per capita. A linha vermelha corresponde ao efeito com o segmento dos $N$ munic\'{i}pios com menor renda per capita,
e a linha azul corresponde ao efeito do segmento composto pelo restante dos munic\'{i}pios.
Pontos azul e vermelho claros s\~{a}o estat\'{i}sticamente significantes ao n\'{i}vel de 10\%,
enquanto os pontos escuros n\~{a}o s\~{a}o.}

\end{center}
\end{figure}

%\begin{figure}[!htb]
%\begin{center}
%\begin{overpic}[width=17cm]{figuras/fig_txprenatal_media_segmentado.png}
%\put(0,40){\rotatebox{90}{M\'{e}dia da renda per capita}}
%\put(50,2){N}
%\end{overpic}
%\caption{\label{fig_txprenatal_media_segmentado} M\'{e}dia da renda per capita.
%A linha vermelha corresponde \`{a} m\'{e}dia dos $N$ munic\'{i}pios com menor renda per capita,
%e a linha azul corresponde \`{a} m\'{e}dia do restante dos munic\'{i}pios.
%A reta tracejada preta corresponde \`{a} m\'{e}dia de todos os munic\'{i}pios, igual a $13.66178$.}
%\end{center}
%\end{figure}


Para $N < 800$, o efeito do tratamento para os $N$ munic\'{i}pios com menor renda per capita n\~{a}o \'{e} estatisticamente significante,
com exce\c{c}\~{a}o de alguns valores, como para $N=400$ ou $N=475$. A partir deste valor,
pode-se perceber que o efeito \'{e} predominantemente maior para os municípios com menor renda.
Cabe notar também que para todo $N$ no intervalo $[2000,4000]$, o efeito dos segmentos \'{e} aproximadamente sim\'{e}trico sobre a m\'{e}dia, assumindo o valor de aproximadamente $1.54$.
Para $N > 4400$, o efeito do tratamento para os munic\'{i}pios restantes n\~{a}o \'{e} estatisticamente significante para a maior parte dos valores.


%------------------------------------------------


\section{Conclus\~{a}o \label{secConclusao}}

O artigo parte do problema apontado pela literatura quanto a ineficácia do Programa Mais Médicos em propiciar melhoria no estado de saúde da população, e busca evidências de seus efeitos ainda no curto prazo. Para isso, faz uso da técnica de estimação por Diferenças em Diferenças, considerando elementos de robustez a partir do controle por covariadas socioeconômicas, ambientais e relativas a serviços médicos. Tudo isso condicionado à existência de adequada correlação entre o estado de saúde da população e a propor\c{c}\~{a}o de nascidos vivos de m\~{a}es com sete ou mais consultas pr\'{e}-natal.

Os resultados estão de acordo com as estratégias empíricas e são favoráveis à existência de efeito de curto prazo. Quanto a isso, a estima\c{c}\~{a}o do modelo estrutural por efeitos fixos, al\'{e}m de eliminar efeitos individuais n\~{a}o-observ\'{a}veis e atenuar problemas de endogeneidade, bastante comum aos munic\'{i}pios brasileiros, sugere que o efeito do programa sobre a variável de interesse é positivo e bastante consistente com os controles. O impacto foi estimado em cerca de 1.8, em média, no modelo básico, controlando por aspectos socioeconsômicos, ambientais e no modelo com todos os controles. Em todos os casos, os resultados contribuem para validar a hipótese de que o programa já produziu efeito sobre a saúde da população em 2014, um ano após sua implantação. 

%As estimativas por diferen\c{c}as em diferen\c{c}as encontradas s\~{a}o mostradas na Tabela \ref{tabEfeitosPMM}.
%Todas s\~{a}o bem pr\'{o}ximas, com ou sem o uso de covariadas, o que mostra uma consist\^{e}ncia no resultado encontrado.
%
%\begin{table}[!htb]\centering \caption{Efeitos do Programa Mais M\'{e}dicos \label{tabEfeitosPMM}}
%\begin{tabular}{|p{12cm}|c|c|}\hline\hline
%\textbf{Modelo} & \textbf{Equa\c{c}\~{a}o} & \textbf{Efeito} \\ \hline
%Diferen\c{c}as em diferen\c{c}as simples & (\ref{eq_mod_txprenatal}) & 1.760 \\ \hline
%Diferen\c{c}as em diferen\c{c}as usando covariadas s\'{o}cio-econ\^{o}micas & (\ref{eq_mod_txprenatal_cov1}) & 1.759 \\ \hline
%Diferen\c{c}as em diferen\c{c}as usando covariadas ambientais & (\ref{eq_mod_txprenatal_cov2}) & 1.831 \\ \hline
%Diferen\c{c}as em diferen\c{c}as usando covariadas s\'{o}cio-econ\^{o}micas, ambientais e relacionadas a servi\c{c}os m\'{e}dicos & (\ref{eq_mod_txprenatal_cov3}) & 1.731 \\ 
%\hline
%\end{tabular}
%\end{table}

Foi considerado também que o impacto do programa deveria ter efeito acentuado nos municípios com nível de renda per capita mais baixo, em consonância com um de seus objetivos de promover a oferta de médicos em regiões com relativa proporção da população em condição de pobreza extrema. De uma forma geral, a regress\~{a}o segmentada por pib per capita mostra que o programa contribui para os munic\'{i}pios mais pobres de uma forma mais significativa que os mais ricos.

Porém, alguns resultados inesperados exigem uma melhor análise. A exemplo, o efeito renda requer uma análise minuciosa por faixas de renda a fim de se entender melhor seu sinal negativo, bem como obter uma estimativa de sua elasticidade. Outros resultados que demandam melhor apreciação são os sinais inesperados de agua tratada e equipamentos.

Outras limitações estão relacionadas a própria base de dados, sobretudo em relação ao método de coleta dos dados do Siab. A presença de viés nos dados poderá suportar sub ou super-estimação dos valores aqui apresentados. Como sugest\~{a}o de trabalhos futuros, recomenda-se realizar um \textit{matching}
para utilizar os munic\'{i}pios no suporte comum, além de verificar o efeito do programa sobre outras vari\'{a}veis de interesse, tais como:
taxa de mortalidade, propor\c{c}\~{a}o de interna\c{c}\~{o}es por condi\c{c}\~{o}es sens\'{i}veis \`{a} Aten\c{c}\~{a}o B\'{a}sica,
e cobertura de acompanhamento das condicionalidades de Sa\'{u}de do Programa Bolsa Fam\'{i}lia.


%----------------------------------------------------------------------------------------
%	REFERENCE LIST
%----------------------------------------------------------------------------------------

% \begin{thebibliography}{99} % Bibliography - this is intentionally simple in this template
% 
% \bibitem[Figueredo and Wolf, 2009]{Figueredo:2009dg}
% Figueredo, A.~J. and Wolf, P. S.~A. (2009).
% \newblock Assortative pairing and life history strategy - a cross-cultural
%   study.
% \newblock {\em Human Nature}, 20:317--330.
%  
% \end{thebibliography}
\newpage
\bibliography{referencias}

\clearpage

%------------------------------------------------

\section{Ap\^{e}ndice \label{secApendice}}

\subsection{Detalhamento das bases de dados utilizadas no artigo \label{secApendice_BDs}}

\begin{table}[!htb]\centering \caption{Bancos de dados \label{tabBDs}}
\begin{tabular}{|p{3cm}|p{5cm}|p{8cm}|}\hline\hline
\textbf{Fornecedor} & \textbf{Banco de Dados} & \textbf{Endere\c{c}o} \\ \hline
INEP & Indicadores Educacionais & \url{http://portal.inep.gov.br/indicadores-educacionais} \\ \hline
DAB - Departamento de Aten\c{c}\~{a}o B\'{a}sica & Hist\'{o}rico de Cobertura da Sa\'{u}de da Fam\'{i}lia & \url{http://dab.saude.gov.br/portaldab/historico_cobertura_sf.php} \\ \hline
\multirow{8}{3cm}{DATASUS - Departamento de Inform\'{a}tica do Sistema \'{U}nico de Sa\'{u}de (SUS)} & Sistema de Informa\c{c}\~{a}o de Aten\c{c}\~{a}o B\'{a}sica (Siab) – Situa\c{c}\~{a}o de Saneamento & \url{http://tabnet.datasus.gov.br/cgi/deftohtm.exe?siab/cnv/SIABCbr.def} \\ \cline{2-3}
 & Sistema de Informa\c{c}\~{a}o de Aten\c{c}\~{a}o B\'{a}sica (Siab) – Produ\c{c}\~{a}o e Marcadores & \url{http://tabnet.datasus.gov.br/cgi/deftohtm.exe?siab/cnv/SIABPbr.def} \\ \cline{2-3}
 & Sistema de Informa\c{c}\~{o}es sobre Or\c{c}amentos P\'{u}blicos em Sa\'{u}de (SIOPS) & \url{http://siops-asp.datasus.gov.br/CGI/deftohtm.exe?SIOPS/serhist/municipio/mIndicadores.def} \\ \cline{2-3}
 & CNES - Recursos F\'{i}sicos - Equipamentos & \url{http://tabnet.datasus.gov.br/cgi/deftohtm.exe?cnes/cnv/equipobr.def} \\ \cline{2-3}
 & CNES - Recursos Humanos - Profissionais - Indiv\'{i}duos & \url{http://tabnet.datasus.gov.br/cgi/deftohtm.exe?cnes/cnv/prid02br.def} \\ \cline{2-3}
 & Indicadores Municipais do rol de Diretrizes, Objetivos, Metas e Indicadores 2015 - Valores Absolutos & \url{http://tabnet.datasus.gov.br/cgi/tabcgi.exe?pacto/2015/cnv/absbr.def} \\ \cline{2-3}
 & Indicadores Municipais do rol de Diretrizes, Objetivos, Metas e Indicadores 2015 & \url{http://tabnet.datasus.gov.br/cgi/deftohtm.exe?pacto/2015/cnv/coapmunbr.def} \\ \cline{2-3}
 & Nascidos Vivos & \url{http://tabnet.datasus.gov.br/cgi/deftohtm.exe?sinasc/cnv/nvbr.def} \\ \hline
\multirow{3}{*}{IBGE} & Populacao residente (2010) & \url{http://www.ibge.gov.br/home/estatistica/populacao/censo2010/indicadores_sociais_municipais/indicadores_sociais_municipais_tab_uf_zip.shtm} \\ \cline{2-3}
 & Estimativas populacionais para os munic\'{i}pios brasileiros enviadas para o TCU & \url{http://www.ibge.gov.br/home/estatistica/populacao/estimativa2014/default.shtm} \\ \cline{2-3}
 & Produto Interno Bruto dos Municipios 2010-2013 & \url{http://www.ibge.gov.br/home/estatistica/economia/pibmunicipios/2010_2013/default.shtm} \\
\hline
\end{tabular}
\begin{tablenotes}
\item \tiny{Fonte: Elaboração dos autores.}
\end{tablenotes}
\end{table}

\clearpage

\subsection{Detalhamento das vari\'{a}veis usadas nas regress\~{o}es \label{secApendice_Variaveis}}

\begin{table}[!htb]\centering \caption{Vari\'{a}veis envolvidas nas regress\~{o}es \label{tabVariaveis}}
\begin{tabular}{|c|p{12cm}|}\hline\hline
\textbf{Vari\'{a}vel} & \textbf{Descri\c{c}\~{a}o} \\ \hline
$PMM$ & Vari\'{a}vel dummy que indica se o munic\'{i}pio j\'{a} recebeu m\'{e}dicos pelo Programa Mais M\'{e}dicos \\ \hline
$pib$ & PIB com o ano de 2014 estimado pela eq. (\ref{eq_estimativa_pib_municipal_1}) \\ \hline
$pibreal$ & PIB real. O deflator utilizado sobre $pib$ foi o IPCA. \\ \hline
$pop$ & Popula\c{c}\~{a}o municipal fornecida para o TCU \\ \hline
$pibrealpc$ & PIB real per capita $\Rightarrow pibrealpc = pibreal / pop$ \\ \hline
$distorcao\_ef$ & Taxa de distor\c{c}\~{a}o idade-s\'{e}rie total no ensino fundamental \\ \hline
$abandono\_em$ & Taxa de abandono total no ensino m\'{e}dio \\ \hline
$aguaclorada$ & Domic\'{i}lios com \'{a}gua clorada \\ \hline
$aguafervida$ & Domic\'{i}lios com \'{a}gua fervida \\ \hline
$aguafiltrada$ & Domic\'{i}lios com \'{a}gua filtrada \\ \hline
$aguapublica$ & Domic\'{i}lios servidos de \'{a}gua proveniente de uma rede geral \\ \hline
$aguatratada$ & $aguatratada = aguaclorada + aguafervida + aguafiltrada + aguapublica$ \\ \hline
$coletalixo$ & Domic\'{i}lios com lixo coletado \\ \hline
$despesasps$ & Despesas com recursos pr\'{o}prios \\ \hline
$despesaspsreal$ & $despesasps$ deflacionada. \\ \hline
$equipamentos$ & Equipamentos existentes em uso pelo SUS tendo como base o m\^{e}s de maio. \\ \hline
$implantados\_esf$ & N\'{u}mero de eSF (Estrat\'{e}gia Sa\'{u}de da Fam\'{i}lia) aptas para o munic\'{i}pio receber incentivos financeiros. \\
\hline
\end{tabular}
\begin{tablenotes}
\item \tiny{Fonte: Elaboração dos autores.}
\end{tablenotes}
\end{table}


% \clearpage
% 
% \subsection{Sa\'{i}das das regress\~{o}es \label{secApendice_saidas}}
% 
% \subsubsection{Sa\'{i}das das regress\~{o}es da Tabela \ref{tab_txprenatal} \label{apendice_saida_tab_txprenatal}}
% 
% O resultado mostrados nas colunas (1), (2) e (3) da Tabela \ref{tab_txprenatal}
% s\~{a}o obtidos por meio das regress\~{o}es das Sa\'{i}das \ref{saida_txprenatal_teste_tendencia},
% \ref{saida_txprenatal_regress} e \ref{saida_txprenatal_xtreg},
% respectivamente.
% 
% \begin{Output}[Sa\'{i}da da coluna (1) da Tabela \ref{tab_txprenatal}.]
% \begin{OutputListing}
% . regress txprenatal PMM d2012 d2014 PMM_d2012 PMM_d2014
% 
%       Source |       SS       df       MS              Number of obs =   16686
% -------------+------------------------------           F(  5, 16680) =  153.38
%        Model |  285395.555     5  57079.1109           Prob > F      =  0.0000
%     Residual |  6207411.65 16680  372.146981           R-squared     =  0.0440
% -------------+------------------------------           Adj R-squared =  0.0437
%        Total |   6492807.2 16685  389.140378           Root MSE      =  19.291
% 
% ------------------------------------------------------------------------------
%   txprenatal |      Coef.   Std. Err.      t    P>|t|     [95% Conf. Interval]
% -------------+----------------------------------------------------------------
%          PMM |  -7.780551    .553682   -14.05   0.000    -8.865827   -6.695276
%        d2012 |    2.48497   .6448307     3.85   0.000     1.221033    3.748907
%        d2014 |   4.849586   .6448307     7.52   0.000      3.58565    6.113523
%    PMM_d2012 |   1.082974   .7830246     1.38   0.167    -.4518373    2.617785
%    PMM_d2014 |   2.843269   .7830246     3.63   0.000     1.308457     4.37808
%        _cons |   66.03249   .4559642   144.82   0.000     65.13875    66.92623
% ------------------------------------------------------------------------------
% \end{OutputListing}
% \label{saida_txprenatal_teste_tendencia}
% \end{Output}
% 
% \begin{Output}[Sa\'{i}da da coluna (2) da Tabela \ref{tab_txprenatal}.]
% \begin{OutputListing}
% . regress txprenatal PMM d2014 PMM_d2014 if (ano==2012) | (ano==2014)
% 
%       Source |       SS       df       MS              Number of obs =   11124
% -------------+------------------------------           F(  3, 11120) =  127.25
%        Model |  119259.249     3  39753.0829           Prob > F      =  0.0000
%     Residual |  3473782.49 11120  312.390512           R-squared     =  0.0332
% -------------+------------------------------           Adj R-squared =  0.0329
%        Total |  3593041.74 11123  323.028117           Root MSE      =  17.675
% 
% ------------------------------------------------------------------------------
%   txprenatal |      Coef.   Std. Err.      t    P>|t|     [95% Conf. Interval]
% -------------+----------------------------------------------------------------
%          PMM |  -6.697577    .507285   -13.20   0.000    -7.691946   -5.703209
%        d2014 |   2.364617   .5907958     4.00   0.000     1.206552    3.522681
%    PMM_d2014 |   1.760295   .7174094     2.45   0.014     .3540451    3.166544
%        _cons |   68.51746   .4177557   164.01   0.000     67.69859    69.33634
% ------------------------------------------------------------------------------
% \end{OutputListing}
% \label{saida_txprenatal_regress}
% \end{Output}
% 
% \begin{Output}[Sa\'{i}da da coluna (3) da Tabela \ref{tab_txprenatal}.]
% \begin{OutputListing}
% . xtset codigo ano, delta(2)
%        panel variable:  codigo (strongly balanced)
%         time variable:  ano, 2010 to 2014
%                 delta:  2 units
% 
% . xtreg txprenatal d2014 PMM_d2014 if (ano==2012) | (ano==2014), fe
% 
% Fixed-effects (within) regression               Number of obs      =     11124
% Group variable: codigo                          Number of groups   =      5562
% 
% R-sq:  within  = 0.1123                         Obs per group: min =         2
%        between = 0.0252                                        avg =       2.0
%        overall = 0.0043                                        max =         2
% 
%                                                 F(2,5560)          =    351.66
% corr(u_i, Xb)  = -0.0400                        Prob > F           =    0.0000
% 
% ------------------------------------------------------------------------------
%   txprenatal |      Coef.   Std. Err.      t    P>|t|     [95% Conf. Interval]
% -------------+----------------------------------------------------------------
%        d2014 |   2.364617   .2427544     9.74   0.000     1.888723     2.84051
%    PMM_d2014 |   1.760295   .2947791     5.97   0.000     1.182412    2.338177
%        _cons |   63.97534   .0973785   656.98   0.000     63.78444    64.16624
% -------------+----------------------------------------------------------------
%      sigma_u |  17.197679
%      sigma_e |  7.2623737
%          rho |  .84866073   (fraction of variance due to u_i)
% ------------------------------------------------------------------------------
% F test that all u_i=0:     F(5561, 5560) =    11.25          Prob > F = 0.0000
% \end{OutputListing}
% \label{saida_txprenatal_xtreg}
% \end{Output}
% 
% %----------------------------------------------------------------------------------------
% \subsubsection{Sa\'{i}das das regress\~{o}es da Tabela \ref{tab_txprenatal_cov1} \label{apendice_saida_tab_txprenatal_cov1}}
% 
% O resultado mostrados nas colunas (1), (2) e (3) da Tabela \ref{tab_txprenatal_cov1}
% s\~{a}o obtidos por meio das regress\~{o}es das Sa\'{i}das \ref{saida_txprenatal_cov1_teste_tendencia},
% \ref{saida_txprenatal_cov1_regress} e \ref{saida_txprenatal_cov1_xtreg},
% respectivamente.
% 
% \begin{Output}[Sa\'{i}da da coluna (1) da Tabela \ref{tab_txprenatal}.]
% \begin{OutputListing}
% regress txprenatal PMM d2012 d2014 PMM_d2012 PMM_d2014 pibrealpc distorcao_ef abandono_em
% 
% 
%       Source |       SS       df       MS              Number of obs =   16646
% -------------+------------------------------           F(  8, 16637) = 1451.55
%        Model |  2659906.99     8  332488.374           Prob > F      =  0.0000
%     Residual |  3810820.84 16637  229.056972           R-squared     =  0.4111
% -------------+------------------------------           Adj R-squared =  0.4108
%        Total |  6470727.83 16645  388.749044           Root MSE      =  15.135
% 
% ------------------------------------------------------------------------------
%   txprenatal |      Coef.   Std. Err.      t    P>|t|     [95% Conf. Interval]
% -------------+----------------------------------------------------------------
%          PMM |  -4.203328   .4367716    -9.62   0.000    -5.059446   -3.347209
%        d2012 |   .2142153   .5079573     0.42   0.673    -.7814352    1.209866
%        d2014 |  -.4201451   .5101791    -0.82   0.410    -1.420151    .5798604
%    PMM_d2012 |     .96979    .615702     1.58   0.115    -.2370516    2.176632
%    PMM_d2014 |   2.683453   .6155324     4.36   0.000     1.476944    3.889962
%    pibrealpc |   .0820234   .0073058    11.23   0.000     .0677032    .0963435
% distorcao_ef |  -.9172876   .0117517   -78.06   0.000    -.9403223    -.894253
%  abandono_em |  -.3213594   .0197085   -16.31   0.000    -.3599901   -.2827286
%        _cons |   89.44527   .4678602   191.18   0.000     88.52821    90.36233
% ------------------------------------------------------------------------------
% \end{OutputListing}
% \label{saida_txprenatal_cov1_teste_tendencia}
% \end{Output}
% 
% \begin{Output}[Sa\'{i}da da coluna (2) da Tabela \ref{tab_txprenatal}.]
% \begin{OutputListing}
% regress txprenatal PMM d2014 PMM_d2014 pibrealpc distorcao_ef abandono_em if (ano==2012) | (ano==2014)
% 
% 
%       Source |       SS       df       MS              Number of obs =   11098
% -------------+------------------------------           F(  6, 11091) = 1118.85
%        Model |  1350066.92     6  225011.154           Prob > F      =  0.0000
%     Residual |  2230505.45 11091  201.109499           R-squared     =  0.3771
% -------------+------------------------------           Adj R-squared =  0.3767
%        Total |  3580572.38 11097  322.661294           Root MSE      =  14.181
% 
% ------------------------------------------------------------------------------
%   txprenatal |      Coef.   Std. Err.      t    P>|t|     [95% Conf. Interval]
% -------------+----------------------------------------------------------------
%          PMM |  -3.581278   .4101392    -8.73   0.000    -4.385224   -2.777332
%        d2014 |  -.3733802   .4771219    -0.78   0.434    -1.308624    .5618636
%    PMM_d2014 |   1.721386   .5768472     2.98   0.003     .5906631    2.852109
%    pibrealpc |   .0563019   .0078985     7.13   0.000     .0408195    .0717843
% distorcao_ef |  -.8365725   .0135098   -61.92   0.000     -.863054   -.8100909
%  abandono_em |  -.3038758   .0225825   -13.46   0.000    -.3481415   -.2596101
%        _cons |   88.18838   .4705375   187.42   0.000     87.26605    89.11072
% ------------------------------------------------------------------------------
% \end{OutputListing}
% \label{saida_txprenatal_cov1_regress}
% \end{Output}
% 
% \begin{Output}[Sa\'{i}da da coluna (3) da Tabela \ref{tab_txprenatal}.]
% \begin{OutputListing}
% xtset codigo ano, delta(2)
%        panel variable:  codigo (strongly balanced)
%         time variable:  ano, 2010 to 2014
%                 delta:  2 units
% 
% xtreg txprenatal d2014 PMM_d2014 pibrealpc distorcao_ef abandono_em if (ano==2012) | (ano==2014), fe
% 
% 
% Fixed-effects (within) regression               Number of obs      =     11098
% Group variable: codigo                          Number of groups   =      5551
% 
% R-sq:  within  = 0.1168                         Obs per group: min =         1
%        between = 0.2630                                        avg =       2.0
%        overall = 0.1899                                        max =         2
% 
%                                                 F(5,5542)          =    146.61
% corr(u_i, Xb)  = 0.3102                         Prob > F           =    0.0000
% 
% ------------------------------------------------------------------------------
%   txprenatal |      Coef.   Std. Err.      t    P>|t|     [95% Conf. Interval]
% -------------+----------------------------------------------------------------
%        d2014 |   1.849812   .2721851     6.80   0.000     1.316222    2.383402
%    PMM_d2014 |   1.759231    .294749     5.97   0.000     1.181408    2.337055
%    pibrealpc |  -.0558441   .0323849    -1.72   0.085    -.1193312    .0076429
% distorcao_ef |  -.1758325   .0439009    -4.01   0.000    -.2618955   -.0897696
%  abandono_em |  -.0632946   .0230678    -2.74   0.006    -.1085165   -.0180727
%        _cons |   69.48946   1.131922    61.39   0.000     67.27045    71.70847
% -------------+----------------------------------------------------------------
%      sigma_u |  16.133452
%      sigma_e |  7.2426332
%          rho |  .83227261   (fraction of variance due to u_i)
% ------------------------------------------------------------------------------
% F test that all u_i=0:     F(5550, 5542) =     6.73          Prob > F = 0.0000
% \end{OutputListing}
% \label{saida_txprenatal_cov1_xtreg}
% \end{Output}
% 
% 
% %----------------------------------------------------------------------------------------
% 
% \subsubsection{Sa\'{i}das das regress\~{o}es da Tabela \ref{tab_txprenatal_cov2} \label{apendice_saida_tab_txprenatal_cov2}}
% 
% O resultado mostrados nas colunas (1), (2) e (3) da Tabela \ref{tab_txprenatal_cov2}
% s\~{a}o obtidos por meio das regress\~{o}es das Sa\'{i}das \ref{saida_txprenatal_cov2_teste_tendencia},
% \ref{saida_txprenatal_cov2_regress} e \ref{saida_txprenatal_cov2_xtreg},
% respectivamente.
% 
% \begin{Output}[Sa\'{i}da da coluna (1) da Tabela \ref{tab_txprenatal}.]
% \begin{OutputListing}
% regress txprenatal PMM d2012 d2014 PMM_d2012 PMM_d2014 log_aguatratada log_coletalixo
% 
% 
%       Source |       SS       df       MS              Number of obs =   16686
% -------------+------------------------------           F(  7, 16678) =  501.86
%        Model |  1129687.23     7  161383.889           Prob > F      =  0.0000
%     Residual |  5363119.98 16678  321.568532           R-squared     =  0.1740
% -------------+------------------------------           Adj R-squared =  0.1736
%        Total |   6492807.2 16685  389.140378           Root MSE      =  17.932
% 
% ---------------------------------------------------------------------------------
%      txprenatal |      Coef.   Std. Err.      t    P>|t|     [95% Conf. Interval]
% ----------------+----------------------------------------------------------------
%             PMM |  -7.413684   .5154238   -14.38   0.000     -8.42397   -6.403399
%           d2012 |   2.156894   .5994527     3.60   0.000     .9819032    3.331885
%           d2014 |   3.839587   .6002316     6.40   0.000     2.663069    5.016104
%       PMM_d2012 |   1.073447   .7278726     1.47   0.140    -.3532603    2.500155
%       PMM_d2014 |   2.716943    .727876     3.73   0.000     1.290229    4.143658
% log_aguatratada |  -6.567874   .1282733   -51.20   0.000    -6.819303   -6.316444
%  log_coletalixo |   6.582908   .1338342    49.19   0.000     6.320579    6.845237
%           _cons |   74.98954   .6400845   117.16   0.000     73.73491    76.24417
% ---------------------------------------------------------------------------------
% \end{OutputListing}
% \label{saida_txprenatal_cov2_teste_tendencia}
% \end{Output}
% 
% \begin{Output}[Sa\'{i}da da coluna (2) da Tabela \ref{tab_txprenatal}.]
% \begin{OutputListing}
% regress txprenatal PMM d2014 PMM_d2014 log_aguatratada log_coletalixo if (ano==2012) | (ano==2014)
% 
% 
%       Source |       SS       df       MS              Number of obs =   11124
% -------------+------------------------------           F(  5, 11118) =  348.57
%        Model |  486916.795     5   97383.359           Prob > F      =  0.0000
%     Residual |  3106124.94 11118  279.378031           R-squared     =  0.1355
% -------------+------------------------------           Adj R-squared =  0.1351
%        Total |  3593041.74 11123  323.028117           Root MSE      =  16.715
% 
% ---------------------------------------------------------------------------------
%      txprenatal |      Coef.   Std. Err.      t    P>|t|     [95% Conf. Interval]
% ----------------+----------------------------------------------------------------
%             PMM |  -6.482291   .4807244   -13.48   0.000    -7.424596   -5.539986
%           d2014 |   1.878306   .5596682     3.36   0.001     .7812574    2.975355
%       PMM_d2014 |   1.666949   .6784507     2.46   0.014     .3370653    2.996833
% log_aguatratada |  -5.468896   .1507695   -36.27   0.000    -5.764431   -5.173361
%  log_coletalixo |   5.577027   .1582122    35.25   0.000     5.266903    5.887151
%           _cons |   74.36364   .6513158   114.17   0.000     73.08695    75.64034
% ---------------------------------------------------------------------------------
% \end{OutputListing}
% \label{saida_txprenatal_cov2_regress}
% \end{Output}
% 
% \begin{Output}[Sa\'{i}da da coluna (3) da Tabela \ref{tab_txprenatal}.]
% \begin{OutputListing}
% xtset codigo ano, delta(2)
%        panel variable:  codigo (strongly balanced)
%         time variable:  ano, 2010 to 2014
%                 delta:  2 units
% 
% xtreg txprenatal d2014 PMM_d2014 log_aguatratada log_coletalixo if (ano==2012) | (ano==2014), fe
% 
% 
% Fixed-effects (within) regression               Number of obs      =     11124
% Group variable: codigo                          Number of groups   =      5562
% 
% R-sq:  within  = 0.1130                         Obs per group: min =         2
%        between = 0.0317                                        avg =       2.0
%        overall = 0.0241                                        max =         2
% 
%                                                 F(4,5558)          =    177.06
% corr(u_i, Xb)  = 0.0490                         Prob > F           =    0.0000
% 
% ---------------------------------------------------------------------------------
%      txprenatal |      Coef.   Std. Err.      t    P>|t|     [95% Conf. Interval]
% ----------------+----------------------------------------------------------------
%           d2014 |   2.338034   .2444585     9.56   0.000       1.8588    2.817268
%       PMM_d2014 |   1.752107   .2947366     5.94   0.000     1.174308    2.329906
% log_aguatratada |  -.5399064   .2522177    -2.14   0.032    -1.034352   -.0454611
%  log_coletalixo |   .5757157   .2731382     2.11   0.035     .0402582    1.111173
%           _cons |   64.20027   .4834764   132.79   0.000     63.25247    65.14807
% ----------------+----------------------------------------------------------------
%         sigma_u |  17.017838
%         sigma_e |  7.2606725
%             rho |  .84600163   (fraction of variance due to u_i)
% ---------------------------------------------------------------------------------
% F test that all u_i=0:     F(5561, 5558) =     9.95          Prob > F = 0.0000
% \end{OutputListing}
% \label{saida_txprenatal_cov2_xtreg}
% \end{Output}
% 
% 
% %----------------------------------------------------------------------------------------
% 
% \subsubsection{Sa\'{i}das das regress\~{o}es da Tabela \ref{tab_txprenatal_cov3} \label{apendice_saida_tab_txprenatal_cov3}}
% 
% O resultado mostrados nas colunas (1), (2) e (3) da Tabela \ref{tab_txprenatal_cov3}
% s\~{a}o obtidos por meio das regress\~{o}es das Sa\'{i}das \ref{saida_txprenatal_cov3_teste_tendencia},
% \ref{saida_txprenatal_cov3_regress} e \ref{saida_txprenatal_cov3_xtreg},
% respectivamente.
% 
% \begin{Output}[Sa\'{i}da da coluna (1) da Tabela \ref{tab_txprenatal}.]
% \begin{OutputListing}
% regress txprenatal PMM d2012 d2014 PMM_d2012 PMM_d2014 pibrealpc abandono_em log_aguatratada log_coletalixo log_despesaspsreal equipamentos implantados_esf
% 
% 
%       Source |       SS       df       MS              Number of obs =   16650
% -------------+------------------------------           F( 12, 16637) =  527.01
%        Model |  1782580.78    12  148548.398           Prob > F      =  0.0000
%     Residual |  4689454.85 16637  281.869018           R-squared     =  0.2754
% -------------+------------------------------           Adj R-squared =  0.2749
%        Total |  6472035.63 16649  388.734196           Root MSE      =  16.789
% 
% ------------------------------------------------------------------------------------
%         txprenatal |      Coef.   Std. Err.      t    P>|t|     [95% Conf. Interval]
% -------------------+----------------------------------------------------------------
%                PMM |  -6.550936   .4848462   -13.51   0.000    -7.501286   -5.600586
%              d2012 |   1.479522   .5631438     2.63   0.009     .3756998    2.583343
%              d2014 |    1.51435   .5658334     2.68   0.007     .4052567    2.623444
%          PMM_d2012 |   1.051255   .6830238     1.54   0.124    -.2875444    2.390054
%          PMM_d2014 |   2.845467   .6827891     4.17   0.000     1.507128    4.183807
%          pibrealpc |   .1154016    .008291    13.92   0.000     .0991503    .1316529
%        abandono_em |   -.825521   .0201579   -40.95   0.000    -.8650326   -.7860094
%    log_aguatratada |  -4.988404   .1273387   -39.17   0.000    -5.238001   -4.738806
%     log_coletalixo |    5.03758   .1341323    37.56   0.000     4.774666    5.300493
% log_despesaspsreal |   .3746013   .0428693     8.74   0.000     .2905728    .4586298
%       equipamentos |    .000328   .0000728     4.50   0.000     .0001853    .0004708
%    implantados_esf |  -.1060787   .0104962   -10.11   0.000    -.1266523   -.0855051
%              _cons |   68.41732   1.339908    51.06   0.000     65.79096    71.04369
% ------------------------------------------------------------------------------------
% \end{OutputListing}
% \label{saida_txprenatal_cov3_teste_tendencia}
% \end{Output}
% 
% \begin{Output}[Sa\'{i}da da coluna (2) da Tabela \ref{tab_txprenatal}.]
% \begin{OutputListing}
% regress txprenatal PMM d2014 PMM_d2014 pibrealpc abandono_em log_aguatratada log_coletalixo log_despesaspsreal equipamentos implantados_esf if (ano==2012) | (ano==2014)
% 
% 
%       Source |       SS       df       MS              Number of obs =   11102
% -------------+------------------------------           F( 10, 11091) =  334.60
%        Model |  830117.837    10  83011.7837           Prob > F      =  0.0000
%     Residual |  2751566.45 11091  248.090024           R-squared     =  0.2318
% -------------+------------------------------           Adj R-squared =  0.2311
%        Total |  3581684.29 11101  322.645193           Root MSE      =  15.751
% 
% ------------------------------------------------------------------------------------
%         txprenatal |      Coef.   Std. Err.      t    P>|t|     [95% Conf. Interval]
% -------------------+----------------------------------------------------------------
%                PMM |   -5.73789   .4551511   -12.61   0.000    -6.630067   -4.845713
%              d2014 |   .3563073   .5309593     0.67   0.502    -.6844673    1.397082
%          PMM_d2014 |   1.796626   .6406196     2.80   0.005     .5408975    3.052354
%          pibrealpc |   .0871109   .0089495     9.73   0.000     .0695683    .1046534
%        abandono_em |  -.7298909   .0234746   -31.09   0.000    -.7759054   -.6838765
%    log_aguatratada |  -4.176314   .1499986   -27.84   0.000    -4.470338    -3.88229
%     log_coletalixo |   4.282168   .1588847    26.95   0.000     3.970725     4.59361
% log_despesaspsreal |   .3819426   .0447246     8.54   0.000     .2942744    .4696109
%       equipamentos |   .0003084   .0000793     3.89   0.000      .000153    .0004638
%    implantados_esf |  -.1057519   .0113295    -9.33   0.000    -.1279598    -.083544
%              _cons |   67.36569   1.401172    48.08   0.000     64.61914    70.11223
% ------------------------------------------------------------------------------------
% \end{OutputListing}
% \label{saida_txprenatal_cov3_regress}
% \end{Output}
% 
% \begin{Output}[Sa\'{i}da da coluna (3) da Tabela \ref{tab_txprenatal}.]
% \begin{OutputListing}
% xtset codigo ano, delta(2)
%        panel variable:  codigo (strongly balanced)
%         time variable:  ano, 2010 to 2014
%                 delta:  2 units
% 
% xtreg txprenatal d2014 PMM_d2014 pibrealpc abandono_em log_aguatratada log_coletalixo log_despesaspsreal equipamentos implantados_esf if (ano==2012) | (ano==2014), fe
% 
% 
% Fixed-effects (within) regression               Number of obs      =     11102
% Group variable: codigo                          Number of groups   =      5552
% 
% R-sq:  within  = 0.1175                         Obs per group: min =         1
%        between = 0.0007                                        avg =       2.0
%        overall = 0.0017                                        max =         2
% 
%                                                 F(9,5541)          =     81.94
% corr(u_i, Xb)  = -0.1376                        Prob > F           =    0.0000
% 
% ------------------------------------------------------------------------------------
%         txprenatal |      Coef.   Std. Err.      t    P>|t|     [95% Conf. Interval]
% -------------------+----------------------------------------------------------------
%              d2014 |   2.210503   .2507264     8.82   0.000      1.71898    2.702025
%          PMM_d2014 |   1.730737   .2958385     5.85   0.000     1.150778    2.310697
%          pibrealpc |  -.0540772   .0323889    -1.67   0.095    -.1175722    .0094178
%        abandono_em |  -.0603519   .0230862    -2.61   0.009    -.1056099    -.015094
%    log_aguatratada |  -.4624339   .2522578    -1.83   0.067    -.9569581    .0320904
%     log_coletalixo |   .4931951   .2731357     1.81   0.071    -.0422581    1.028648
% log_despesaspsreal |   .0921033   .0330894     2.78   0.005     .0272351    .1569716
%       equipamentos |  -.0011803   .0004825    -2.45   0.014    -.0021263   -.0002344
%    implantados_esf |   .1072242    .043713     2.45   0.014     .0215296    .1929187
%              _cons |   62.47469   1.221156    51.16   0.000     60.08075    64.86864
% -------------------+----------------------------------------------------------------
%            sigma_u |   17.36878
%            sigma_e |  7.2428409
%                rho |  .85186722   (fraction of variance due to u_i)
% ------------------------------------------------------------------------------------
% F test that all u_i=0:     F(5551, 5541) =     8.66          Prob > F = 0.0000
% \end{OutputListing}
% \label{saida_txprenatal_cov3_xtreg}
% \end{Output}

%----------------------------------------------------------------------------------------


%----------------------------------------------------------------------------------------

\end{document}
